% !TeX root = ../../Book2.tex
% 纲要标题：## 第18章　Brauer 群与循环代数
% 纲要位置：计划纲要.md，第 1626 行。

\chapter{Brauer 群与循环代数}
\label{b2:ch:18}
\BookChapterNumber{b2:ch:18}

\begin{chapterabstract}
矩阵大小改变中心单代数的维数，却不改变其除代数代表。Brauer 群把这部分稳定信息组织成张量积下的交换群。循环代数使其中的一些类能够由一个 Galois 自同构与一个范数参数计算；约化迹和约化范数则把分裂矩阵中的运算准确地降回底域。对有限 Galois 扩张 \(K/k\)，交叉积与\(2\)-余圈描述相对 Brauer 群 \(\operatorname{Br}(K/k)\)；最后证明周期整除指数，并给出周期二、指数四的具体除代数。
\end{chapterabstract}

\begin{learninggoals}
\begin{itemize}
\item 用除代数代表判断 Brauer 等价，区分代数的次数、指数与 Brauer 类的周期。
\item 写出循环代数的乘法及分裂矩阵，把分裂性转化为范数方程。
\item 计算约化特征多项式、迹和范数，说明它们与正则作用及子域运算的关系。
\item 从半线性矩阵记录余圈，核对换基、余边及反代数的方向。
\item 运用有限分裂域次数的整除性，辨认周期与指数可能出现的差别。
\end{itemize}
\end{learninggoals}

若 \(D\) 是域 \(k\) 上的中心除代数，\(D\) 与 \(M_r(D)\) 的维数不同，二者的模却保留了相同的除代数部分。把矩阵大小当作分类数据，便会把这种稳定信息遮住；Brauer 等价正是将矩阵因子略去后比较中心单代数。张量积给出这些类的加法，但要证明它确实形成群，还须给出逆元并核对选择矩阵代表时运算不变。

\ProseParagraphBreak

仅有抽象群运算还难以计算具体类。若一个循环 Galois 扩张已经使代数分裂，就可以把自同构和一个底域元素写进乘法关系，得到循环代数；它何时已经是矩阵代数，又归结到范数方程。另一方面，分裂后矩阵的迹与行列式若要回到底域，必须证明它们不依赖分裂同构。这些可算对象把稳定分类与原代数中的运算连接起来。

\ProseParagraphBreak

本章需要第十七章的中心单代数、可分分裂域、中心化子及 Skolem--Noether 定理，也需要第一卷的有限 Galois 扩张与域范数。\secref{b2:sec:1801}建立 Brauer 等价与群运算，\secref{b2:sec:1802}给出循环代数的直接构造，\secref{b2:sec:1803}处理分裂模型中迹与行列式的下降，\secref{b2:sec:1804}由余圈方程研究相对 Brauer 群，\secref{b2:sec:1805}比较周期与指数。

\ProseParagraphBreak
\(\mathbb C/\mathbb R\)、有限域及二次有理扩张上的循环代数，为一般交叉积提供了可计算的例子。中心单代数、循环代数、约化运算与周期和指数的进一步论述，可对照 Gille–Szamuely 的《Central Simple Algebras and Galois Cohomology》\cite[Sections 2.4--2.8; Chapter 3]{GilleSzamuely2017CentralSimpleAlgebras}。

% !TeX root = ../../Book2.tex
% 纲要标题：### 18.1 Brauer 等价
% 纲要位置：计划纲要.md，第 1628 行。

\section{Brauer 等价}
\label{b2:sec:1801}

第十六章的矩阵环 Morita 等价给出 \(D\text{-Mod}\simeq M_r(D)\text{-Mod}\)，而第十七章把每个中心单代数写成 \(M_r(D)\)。矩阵大小可以改变代数的维数，却不改变它背后的除代数。固定中心域 \(k\) 及其标量作用，Brauer 等价正是记录这种模范畴的稳定类型；张量积将在这些类上给出可以取逆的运算。本章固定底域 \(k\)，所有中心单代数有限维、非零，代数同构均保持 \(k\) 标量。

% 纲要要点（供目录核对）：
%
% * 中心单代数的 Brauer 等价；以固定底域的线性 Morita 等价解释稳定分类，正式证明非零角代数保持 Brauer 类
% * 张量积给出的群结构；核对加矩阵块的相容性，以包络同构解释反代数真正给出逆元
% * 除代数代表、指数与周期的区分；指数等于最小有限分裂次数，正式证明扩域后指数的两项整除约束
% * 利用反代数解释第一类对合使 Brauer 类的二倍为零，区分周期二与指数二，保留次数、指数及周期各自的意义
%

\subsection{除代数代表与 Brauer 等价}
\label{b2:subsec:1801:01}

若 \(A=M_p(D)\)、\(B=M_q(D)\)，希望略去原来的矩阵重数，就应允许两侧再加不同大小的矩阵块，因为
\[
M_q(A)\cong M_{pq}(D)\cong M_p(B).
\]
要求两侧都加同一个矩阵大小，会保留原来的重数差异，不能表达这种稳定关系。

\begin{definition}\label{b2:def:brauer-equivalence}
设 \(k\) 为域，\(A,B\) 为中心单 \(k\) 代数。若存在正整数 \(r,s\)，使
\[
M_r(A)\cong M_s(B)
\]
为 \(k\) 代数同构，则称 \(A,B\) \term[Brauer-dengjia]{Brauer 等价}[Brauer equivalent]。
\end{definition}

\begin{proposition}\label{b2:prop:brauer-division-representative}
设 \(k\) 为域。每个中心单 \(k\) 代数都同构于某个 \(M_r(D)\)，其中 \(r\geq1\)，\(D\) 为中心是 \(k\) 的有限维除代数。若 \(A\cong M_r(D)\)、\(B\cong M_s(E)\) 为两个中心单 \(k\) 代数，则 \(A,B\) 在 Brauer 意义下等价，当且仅当 \(D\cong E\) 为 \(k\) 代数同构。上述条件也等价于幺左模范畴 \(A\text{-Mod}\) 与 \(B\text{-Mod}\) 之间存在 \(k\) 线性等价，即该等价在每个 Hom 空间上诱导的映射均为 \(k\) 线性。因此 Brauer 等价是等价关系，每个等价类具有唯一到 \(k\) 代数同构的除代数代表。对任意非零幂等元 \(e\in A\)，角代数 \(eAe\) 以 \(e\) 为单位，也是中心单 \(k\) 代数，并与 \(A\) 在 Brauer 意义下等价；其 \(k\) 标量为 \(ce\)、\(c\in k\)。
\end{proposition}
\begin{proof}
中心单代数是有限维单 Artin 环，由\cref{b2:prop:simple-artinian}可写成 \(M_r(D)\)。矩阵环的中心恰为 \(Z(D)\) 中元素组成的标量矩阵，故 \(Z(D)=k\)。这个结构同构保持中心中的标量，故为 \(k\) 代数同构。若 \(A=M_r(D),B=M_s(D)\)，则 \(M_s(A)\cong M_{rs}(D)\cong M_r(B)\)，所以它们 Brauer 等价。

反过来，定义中的矩阵同构给出同一个单环的两种矩阵分解。\cref{b2:thm:artin-wedderburn}及其唯一性说明除环代表同构。该同构可由一个简单模的自同态反环恢复；矩阵同构保持 \(k\) 的中心作用，所以恢复的同构也保持 \(k\)。于是上述充要条件成立，自反、对称和传递均随之得到。

矩阵 Morita 等价由 \(k\) 双模的张量积构造，保持 Hom 空间的 \(k\) 标量，所以同一个 \(D\) 给出上述 \(k\) 线性范畴等价。反过来，等价保持简单对象，并给出其自同态环的 \(k\) 代数同构。两侧各只有一种简单模类型，自同态环分别为 \(D^{\mathrm{op}}\)、\(E^{\mathrm{op}}\)，故 \(D\cong E\) 也保持 \(k\) 标量。

最后，在 \(A=M_r(D)=\operatorname{End}_D(D^r)\) 中，幂等算子 \(e\) 给出右 \(D\) 空间分解 \(D^r=\operatorname{im}e\oplus\ker e\)。取适应该分解的基，将 \(e\) 写成 \(\operatorname{diag}(I_s,0)\)、\(s\ge1\)。限制到像空间给出 \(eAe\cong\operatorname{End}_D(\operatorname{im}e)\cong M_s(D)\)，保持 \(k\) 作用和单位。因此角代数中心单，且仍有除代数代表 \(D\)。
\end{proof}

因此每个 Brauer 类可以用唯一的中心除代数同构类型表示。例如，所有 \(M_n(k)\) 都与 \(k\) 在 Brauer 意义下等价。实四元数除代数 \(\mathbb H\) 却不与 \(\mathbb R\) 等价，因为这两个除代数不同构。Brauer 等价不表示原代数同构：矩阵大小仍然决定实际的维数。

\subsection{张量积与 Brauer 群结构}
\label{b2:subsec:1801:02}

\begin{theorem}\label{b2:thm:brauer-group}
设 \(k\) 为域。中心单 \(k\) 代数的 Brauer 等价类\footnote{每个有限维代数可由某个 \(k^n\) 上的有限乘法表表示，全部这些乘法表组成集合，取同构及稳定等价类后仍为集合。}组成交换群，运算、单位及逆分别为
\[
[A]+[B]=[A\otimes_k B],\qquad 0=[k],\qquad -[A]=[A^{\mathrm{op}}].
\]
这个群称为 \(k\) 的\term[Brauer-qun]{Brauer 群}[Brauer group]，记为 \(\operatorname{Br}(k)\)。
\end{theorem}
\symidx{\(\operatorname{Br}(k)\)：域 \(k\) 的 Brauer 群}
\begin{proof}
要使张量积下降到已经略去矩阵大小的类上，须检查加矩阵块与张量积相容。张量积仍中心单，由\cref{b2:thm:csa-tensor-product}已知。对矩阵单位逐项相乘，得到
\[
M_r(A)\otimes_k M_s(B)\cong M_{rs}(A\otimes_k B).
\]
其映射把 \((E_{ij}a)\otimes(E_{uv}b)\) 送到以 \((i,u),(j,v)\) 为行列指标、条目为 \(a\otimes b\) 的矩阵单位。它保持乘法并把相应的向量空间基一一对应，故为同构。这说明张量运算不依赖稳定代表元的选择。

张量积的结合、交换因子同构与单位同构给出交换群运算所需的结合律、交换律和单位。\cref{b2:thm:csa-enveloping-isomorphism}给出
\[
A\otimes_k A^{\mathrm{op}}\cong\operatorname{End}_k(A),
\]
右边为 \(M_{\dim_kA}(k)\)，代表零类，所以反代数确实通过张量积将 \(A\) 抵消为分裂代数。这正是它给出群逆元的原因。
\end{proof}

若 \(K/k\) 为域扩张，标量扩张保持中心单性，并与张量积相容，所以诱导群同态
\[
\operatorname{Br}(k)\longrightarrow\operatorname{Br}(K),\qquad
[A]\longmapsto[A\otimes_kK].
\]
其核记为 \(\operatorname{Br}(K/k)\)，称为\term[xiangdui-Brauer-qun]{相对 Brauer 群}[relative Brauer group]；其中的类恰由被 \(K\) 分裂的代数代表，因为零类等价于除代数代表为底域，亦即代数本身为全矩阵代数。
\symidx{\(\operatorname{Br}(K/k)\)：由 \(K\) 分裂的 Brauer 类}

\ProseParagraphBreak
任意有限域的 Brauer 群为零。事实上，有限维除代数在有限底域上只有有限多个元素，由\cref{b2:thm:wedderburn-little}为交换域；若中心正是底域，它只能等于底域。代数闭域的 Brauer 群也为零，因为每个中心单代数都已分裂。

\subsection{除代数代表、指数与周期}
\label{b2:subsec:1801:03}

\begin{definition}\label{b2:def:brauer-index-period}
设 \(k\) 为域，\(A\) 为中心单 \(k\) 代数，写 \(A\cong M_r(D)\)，其中 \(r\geq1\)，\(D\) 为中心除 \(k\) 代数。记 \(\dim_kD=d^2\)，定义
\[
\operatorname{ind}(A)=d
\]
并称它为 \(A\) 的\term[zhongxindan-daishu-zhishu]{指数}[index of a central simple algebra]。把 Brauer 类 \([A]\in\operatorname{Br}(k)\) 的加法阶称为 \(A\) 的\term[Brauer-zhouqi]{周期}[period]，记为 \(\operatorname{per}(A)\)；此处按群元素阶的通常约定，暂允许其取值为 \(\infty\)。\cref{b2:thm:period-divides-index}将证明周期总是有限的，且整除指数。这里 \(\deg A=\sqrt{\dim_k A}=rd\)，零 Brauer 类的周期和指数均为一。
\end{definition}
\symidx{\(\operatorname{ind}(A)\)：中心单代数的指数}
\symidx{\(\operatorname{per}(A)\)：中心单代数的周期}

指数和周期都不随矩阵大小改变，次数却随之改变。例如 \(M_3(\mathbb H)\) 的次数为六、指数为二。\eqref{b2:eq:quaternion-conjugation}中的四元数共轭给出 \(\mathbb H\cong\mathbb H^{\mathrm{op}}\)，所以其 Brauer 类的二倍为零；它又不是零类，故周期恰为二。这个例子区分了次数与另两个数，但不表示一般代数的周期都等于指数。

\ProseParagraphBreak
设 \(k\) 为任意域，\(A\) 为中心单 \(k\) 代数。若一个 \(k\) 线性映射 \(\sigma:A\to A\) 满足 \(\sigma(1)=1\)、\(\sigma(xy)=\sigma(y)\sigma(x)\) 及 \(\sigma^2=\operatorname{id}\)，则称它为\term[diyilei-duihe]{第一类对合}[involution of the first kind]。它作为 \(A\to A^{\mathrm{op}}\) 的 \(k\) 代数同构，给出 \([A]=[A^{\mathrm{op}}]=-[A]\)，故 \(2[A]=0\)。这个条件限制的是周期，不能据此断言指数至多为二；\cref{b2:exm:period-two-index-four}给出周期二、指数四的例子。

\begin{proposition}\label{b2:prop:index-splitting-degree}
设 \(k\) 为域，\(A\) 为中心单 \(k\) 代数。若有限域扩张 \(L/k\) 分裂 \(A\)，则 \(\operatorname{ind}(A)\mid[L:k]\)。而且存在次数恰为 \(\operatorname{ind}(A)\) 的有限可分分裂域。因此
\[
\operatorname{ind}(A)
=\min\{\,[L:k]\mid L/k\;\text{为分裂}\;A\;\text{的有限域扩张}\,\}.
\]
\end{proposition}
\begin{proof}
令 \(A=M_r(D)\)、\(\deg D=d\)。\(A_L=M_r(D_L)\) 分裂，当且仅当 \(D_L\) 分裂，这由单环矩阵分解的除代数唯一性得到。此时 \(D\) 通过 \(D_L\cong M_d(L)\) 作用在列空间 \(L^d\) 上。把它看成左 \(D\) 向量空间，设维数为 \(t\)。比较 \(k\) 维数得
\[
d[L:k]=t\dim_kD=td^2,
\]
所以 \(d\mid[L:k]\)。\cref{b2:thm:csa-separable-splitting}还为除代数 \(D\) 构造了一个次数为 \(d\) 的可分子域并证明它分裂 \(D\)，同一个域也分裂 \(A\)，得到存在性。
\end{proof}

这说明指数记录最小可能的有限分裂域次数；周期记录在 Brauer 群中反复张量多少次才能得到零类。一个是分裂扩张的大小，另一个是群运算的阶数，定义上不应混同。


\begin{proposition}\label{b2:prop:index-after-base-change}
设 \(k\) 为域，\(A\) 为中心单 \(k\) 代数，\(L/k\) 为域扩张，令 \(d=\operatorname{ind}(A)\)、\(e=\operatorname{ind}(A_L)\)，其中 \(A_L=A\otimes_kL\)。则 \(e\mid d\)。若 \(L/k\) 有限，则还有
\[
\dfrac de\mid[L:k].
\]
特别地，若 \([L:k]\) 与 \(d\) 互素，则 \(e=d\)。
\end{proposition}
\begin{proof}
以除代数代表 \(D\) 代替 \(A\) 不改变扩域后的 Brauer 类。写 \(D_L\cong M_s(E)\)，其中 \(E\) 为中心除 \(L\) 代数，次数为 \(e\)。比较次数得 \(d=se\)，故 \(e\mid d\)。

若 \(L/k\) 有限，通过含幺嵌入 \(D\hookrightarrow D_L\cong M_s(E)\)，列空间 \(E^s\) 成为左 \(D\) 向量空间。记其 \(D\) 维数为 \(t\)，比较底层 \(k\) 维数，得到
\[
s e^2[L:k]=t d^2=t s^2e^2.
\]
这些是正整数之间的等式，消去 \(se^2\) 得 \([L:k]=ts\)，即 \(d/e=s\) 整除扩域次数。若该次数与 \(d\) 互素，\(s\) 同时整除二者，只能等于一，所以 \(e=d\)。
\end{proof}

% !TeX root = ../../Book2.tex
% 纲要标题：### 18.2 循环代数
% 纲要位置：计划纲要.md，第 1637 行。

\section{循环代数}
\label{b2:sec:1802}

一个循环 Galois 扩张不仅能提供标量，还能提供乘法的扭转规则。把扩张的生成自同构写成一个新元素的共轭作用，再规定该元素的某次幂，就得到循环代数。它是否分裂，最后化为原扩张中的一个范数方程。

% 纲要要点（供目录核对）：
%
% * 循环 Galois 扩张及扭曲乘法；先由生成关系产生进位公式，再以两次进位恒等式证明结合律，回扣代数展示的泛性质
% * 循环代数的中心单性；分别核对中心与最少非零系数证明，分裂映射使用原代数的右子域空间
% * 范数与分裂条件；一维子域空间的半线性算子产生范数方程，换生成元直接计算参数及恢复原生成元
%

\subsection{循环 Galois 扩张及扭曲乘法}
\label{b2:subsec:1802:01}

设 \(K/k\) 为次数 \(n\ge2\) 的循环 Galois 扩张，\(\operatorname{Gal}(K/k)=\langle\sigma\rangle\)，并取 \(a\in k^\times\)。希望在包含 \(K\) 的代数中加入一个元素 \(u\)，要求 \(ub=\sigma(b)u\)（\(b\in K\)）及 \(u^n=a\)。前一关系反复使用便给出 \(u^ic=\sigma^i(c)u^i\)；若 \(i+j=qn+r\)、\(0\le r<n\)，后一关系将 \(u^{i+j}\) 化成 \(a^qu^r\)。这两个规则决定了乘积的系数和指数。据此，在左 \(K\) 向量空间
\[
C=K\oplus Ku\oplus\cdots\oplus Ku^{n-1}
\]
上，用下式定义乘法，并按 \(k\) 双线性扩张：
\begin{equation}\label{b2:eq:cyclic-algebra-product}
(bu^i)(cu^j)=b\sigma^i(c)a^{\left\lfloor(i+j)/n\right\rfloor}u^{\,i+j-n\left\lfloor(i+j)/n\right\rfloor},
\qquad 0\le i,j<n.
\end{equation}
这里 \(q=\lfloor(i+j)/n\rfloor\) 记录进位次数，每次进位产生一个 \(a\)。乘法一般不是 \(K\) 双线性的，但 \(\sigma\) 固定 \(k\)，所以仍为 \(k\) 双线性。

\begin{proposition}\label{b2:prop:cyclic-algebra-construction}
设 \(K/k\) 为次数 \(n\geq2\) 的循环 Galois 扩张，\(\operatorname{Gal}(K/k)=\langle\sigma\rangle\)，\(a\in k^\times\)。在左 \(K\) 空间 \(C=\bigoplus_{i=0}^{n-1}Ku^i\) 上，按 \(k\) 双线性延拓乘法
\[
(bu^i)(cu^j)=b\sigma^i(c)a^{\lfloor(i+j)/n\rfloor}
u^{\,i+j-n\lfloor(i+j)/n\rfloor}\qquad(b,c\in K,\ 0\leq i,j<n).
\]
则 \(C\) 是维数为 \(n^2\) 的结合含幺 \(k\) 代数。它由子域 \(K\) 和一个可逆元 \(u\) 生成，满足
\[
ub=\sigma(b)u\quad(b\in K),\qquad u^n=a.
\]
更确切地，给定含幺 \(k\) 代数 \(B\)、指定的含幺 \(k\) 代数嵌入 \(\iota:K\hookrightarrow B\)，以及 \(v\in B\) 满足
\[
v\iota(b)=\iota\bigl(\sigma(b)\bigr)v\quad(b\in K),\qquad v^n=\iota(a),
\]
则存在唯一含幺 \(k\) 代数同态 \(\Phi:C\to B\)，使 \(\Phi|_K=\iota\) 且 \(\Phi(u)=v\)。
\end{proposition}
\begin{proof}
比较三个基型元素 \(bu^i,cu^j,du^\ell\) 的乘积，令 \(q_{i,j}=\lfloor(i+j)/n\rfloor\)、\(r_{i,j}=i+j-nq_{i,j}\)。两种结合次序先分别将 \(i+j\) 或 \(j+\ell\) 化到 \(0,\ldots,n-1\)，总进位数满足
\[
q_{i,j}+q_{r_{i,j},\ell}
=\left\lfloor\frac{i+j+\ell}{n}\right\rfloor
=q_{j,\ell}+q_{i,r_{j,\ell}}.
\]
这是将 \(i+j=nq_{i,j}+r_{i,j}\) 或 \(j+\ell=nq_{j,\ell}+r_{j,\ell}\) 代入中间式所得的恒等式。两种乘积因而含有相同的 \(a\) 次幂，剩下的 \(u\) 指数也都是 \(i+j+\ell-n\lfloor(i+j+\ell)/n\rfloor\)。其余系数都为 \(b\sigma^i(c)\sigma^{i+j}(d)\)，因为 \(\sigma^n=1\)，且 \(a\in k\) 被 \(\sigma\) 固定。因此两种乘积相同，双线性再给出全部元素的结合律。\(1\in K\) 为单位，\(k\) 标量位于中心。底向量空间事先已指定为 \(n\) 个 \(K\) 的直和，所以维数确为 \(n^2\)，无需再假设这些形式单项式无关。

关系由定义直接得到，\(u^{-1}=a^{-1}u^{n-1}\)。给定 \(\iota\) 及 \(v\)，唯一可能的映射把 \(\sum_i b_i u^i\) 送到 \(\sum_i\iota(b_i)v^i\)。由 \(v^i\iota(c)=\iota(\sigma^i(c))v^i\) 与 \(v^n=\iota(a)\)，它保持\eqref{b2:eq:cyclic-algebra-product}，所以确为代数同态。
\end{proof}

该代数称为由这些数据给出的\term[xunhuan-daishu]{循环代数}[cyclic algebra]，记为 \((K/k,\sigma,a)\)。上述泛性质是\cref{b2:prop:algebra-relations-universal}中代数展示原则的具体应用：\(K\) 中的加法与乘法已由 \(\iota\) 保持，新生成元的像只需满足所列的两条扭曲关系，就能唯一确定整个代数同态。
\symidx{\((K/k,\sigma,a)\)：循环扩张及参数确定的循环代数}
当 \(n=1\) 时，另约定 \((k/k,\operatorname{id},a)=k\)；这个退化情形不另引入形式生成元。

\subsection{循环代数的中心单性}
\label{b2:subsec:1802:02}

\begin{theorem}\label{b2:thm:cyclic-algebra-central-simple}
设 \(K/k\) 为次数 \(n\geq1\) 的循环 Galois 扩张，\(\operatorname{Gal}(K/k)=\langle\sigma\rangle\)，\(a\in k^\times\)，则循环代数 \(C=(K/k,\sigma,a)\) 是次数 \(n\) 的中心单 \(k\) 代数，且由 \(K\) 分裂。
\end{theorem}
\begin{proof}
当 \(n=1\) 时 \(C=k\)，结论直接成立。以下设 \(n\ge2\)。

要确定中心，先找出哪些元素与子域 \(K\) 交换，再检查它们是否与生成元 \(u\) 交换。设 \(z=\sum_i b_i u^i\) 与 \(K\) 中所有元素交换。比较 \(zc\) 与 \(cz\)，得 \(b_i\bigl(\sigma^i(c)-c\bigr)=0\) 对所有 \(c\in K\) 成立。对 \(i\ne0\)，\(\sigma^i\ne1\)，可以选 \(c\) 使差非零，故 \(b_i=0\)。所以 \(C\) 中 \(K\) 的中心化子就是 \(K\)。若 \(z=b_0\) 再与 \(u\) 交换，则 \(\sigma(b_0)=b_0\)，故 \(z\in k\)。于是中心恰为 \(k\)。

为证单性，需要证明每个非零双边理想都含有 \(1\)。取这样的理想 \(I\)，并在其中选取非零系数个数最少的元素。右乘 \(u^r\)（\(0\le r<n\)）使指数 \(i\) 变为 \((i+r)\bmod n\)，只是循环搬移各分量，并将进位项的系数乘以非零标量 \(a\)。因此可以将任意一个非零系数搬到常数位置，且非零系数个数不变。再左乘这个常数系数在 \(K\) 中的逆，只将所有系数同时乘以非零数，便得到 \(x=1+\sum_{i>0}b_i u^i\in I\)，仍有最少的非零系数。对 \(c\in K\)，交换子 \(xc-cx\in I\) 的常数项为零，其余非零项只能出现在 \(x\) 原有的非零位置，故非零项数严格小于 \(x\)。若它非零，与极小性矛盾；故它对所有 \(c\) 都为零。由前一段的中心化子计算，\(x\in K\)，进而 \(x=1\)。因此 \(I=C\)，代数单。

构造给出 \(\dim_kC=n^2\)，所以 \([K:k]=n=\deg C\)，恰满足\cref{b2:prop:maximal-field-splits-csa}中子域次数等于代数次数的条件，故 \(K\) 分裂 \(C\)。该命题的实际映射为 \(C\otimes_kK\to\operatorname{End}_K(C)\)，\(c\otimes b\) 作用为 \(x\mapsto cxb\)；目标中的空间是原代数 \(C\)，并未换成 \(C_K\)。把它看成右 \(K\) 空间，其维数为 \(n\)，所以目标是 \(M_n(K)\)。
\end{proof}

在次数二、特征不为二时，若 \(K=k(\sqrt d)\)，则 \(u\sqrt d=-\sqrt d\,u\)，\(u^2=a\)。因而循环代数就是四元数代数 \((d,a)_k\)。这里指定自同构是重要数据；较高次数时，换一个生成元需同时跟踪相应参数，不能只把记号中的 \(\sigma\) 删去。

\subsection{范数与分裂条件}
\label{b2:subsec:1802:03}

循环扩张的范数为
\[
\operatorname{N}_{K/k}(b)=\prod_{j=0}^{n-1}\sigma^j(b).
\]
这是第一卷定理18.2.3的共轭元公式在循环 Galois 扩张上的形式。它给出判断 \(C\) 是否为矩阵代数的条件。

\begin{theorem}[循环代数的分裂判据]\label{b2:thm:cyclic-algebra-norm-splitting}
设 \(K/k\) 为有限循环 Galois 扩张，\(\operatorname{Gal}(K/k)=\langle\sigma\rangle\)，\(a\in k^\times\)，则循环代数 \((K/k,\sigma,a)\) 分裂，当且仅当 \(a\in\operatorname{N}_{K/k}(K^\times)\)。
\end{theorem}
\begin{proof}
当 \([K:k]=1\) 时循环代数为 \(k\)，而域范数为恒等映射，结论成立。以下设 \(n=[K:k]\ge2\)。

若 \(a=\operatorname{N}_{K/k}(b)\)，其中 \(b\in K^\times\)，要构造分裂表示，可以让 \(K\) 在自身上左乘，再寻找与扭曲关系相容的算子 \(T\) 来表示 \(u\)。由 \(uc=\sigma(c)u\) 所要求的相容条件是 \(T(cz)=\sigma(c)T(z)\)。\(K\) 是一维 \(K\) 空间，因此任何满足此条件的算子都由 \(T(1)\) 决定，具体为 \(T(z)=\sigma(z)T(1)\)。取 \(T(1)=b\)，就得到 \(k\) 线性算子 \(T(z)=b\sigma(z)\)，并且
\[
T^n(z)=b\sigma(b)\cdots\sigma^{n-1}(b)\sigma^n(z)
=\operatorname{N}_{K/k}(b)z=az.
\]
因此\cref{b2:prop:cyclic-algebra-construction}的泛性质给出含幺 \(k\) 代数同态 \(C\to\operatorname{End}_k(K)\)，也就是\cref{b2:prop:algebra-modules-representations}所说的 \(C\) 在这个空间上的作用。来源单而映射非零，所以单射；两端维数都是 \(n^2\)，故为同构。

反过来，设 \(C\cong\operatorname{End}_k(V)\)，其中 \(\dim_kV=n\)。这个分裂模型给出 \(C\) 的作用，限制到子域 \(K\) 使 \(V\) 成为左 \(K\) 向量空间。比较 \(\dim_kV=[K:k]\dim_KV\)，得到 \(\dim_KV=1\)。选一个 \(K\) 基向量把 \(V\) 识别为 \(K\)。此时 \(u\) 的作用 \(T\) 必须满足 \(T(cz)=\sigma(c)T(z)\)，所以仍由 \(b=T(1)\) 唯一决定，\(T(z)=b\sigma(z)\)。\(u\) 可逆使 \(b\ne0\)，而 \(u^n=a\) 给出 \(\operatorname{N}_{K/k}(b)=a\)：分裂模型中的一维 \(K\) 空间强制产生了所需的范数解。
\end{proof}

当 \(n\ge2\)、\(b\in K^\times\) 时，把生成元改为 \(v=bu\)。每次将 \(u\) 移过下一个 \(b\) 都施加一次 \(\sigma\)，所以
\[
v^n=b\sigma(b)\cdots\sigma^{n-1}(b)u^n
=\operatorname{N}_{K/k}(b)a.
\]
并且 \(vc=\sigma(c)v\) 对每个 \(c\in K\) 成立。新生成元还能恢复 \(u=b^{-1}v\)，故由构造的泛性质得到从参数为 \(\operatorname{N}_{K/k}(b)a\) 的循环代数到 \(C\) 的满同态；两侧维数都为 \(n^2\)，它是同构。因此参数乘上一个范数不改变代数的同构类型。

\ProseParagraphBreak
例如，取 \(\sigma\) 为复共轭，\(\mathbb C/\mathbb R\) 的范数为 \(z\overline z\)，非零范数恰为正实数。故 \((\mathbb C/\mathbb R,\sigma,a)\) 在 \(a>0\) 时分裂，在 \(a<0\) 时同构于 \(\mathbb H\)。前一种情况可以缩放 \(u\) 使平方为一，后一种可以缩放使平方为负一；范数判据同时解释为什么两种情况不能同构。

% !TeX root = ../../Book2.tex
% 纲要标题：### 18.3 约化迹与约化范数
% 纲要位置：计划纲要.md，第 1643 行。

\section{约化迹与约化范数}
\label{b2:sec:1803}

把中心单代数在一个分裂域中写成矩阵以后，可以取通常的迹与行列式。问题是这些数是否仍属于原来的底域，以及换一个分裂域、换一个矩阵同构后是否改变。本节先解决这两个问题，再把所得运算与域迹、域范数比较。始终令 \(\deg A=n\)。本章通过自然交换同构 \(K\otimes_kA\cong A\otimes_kK\) 采用右侧的标量扩张模型，仍将它记为 \(A_K\)。

% 纲要要点（供目录核对）：
%
% * 在分裂域中的定义；解释抵消半线性、Galois 降系数及任意分裂域比较，Cayley–Hamilton 的系数在中心
% * 与域迹、域范数的比较；同一左乘算子的两种坐标计算及首一因子重数比较，区别四元数标量的最小多项式与约化特征多项式
% * 定义对分裂域选择的独立性与标量扩张相容性；泛型坐标处理扩域后的新元素，约化迹可加、约化范数乘法而不加法，正则迹不能在坏特征除以次数
% * 正式证明一般约化迹配对对称非退化及迹映射满射；纯不可分子域的限制配对可以全退化；对合与转置复合证明保持约化特征多项式和非退化对称迹型
%

\subsection{约化迹与约化范数的定义}
\label{b2:subsec:1803:01}

由\cref{b2:thm:csa-separable-splitting}，可取一个固定的有限 Galois 分裂域 \(K_0/k\) 和同构 \(\rho:A\otimes_kK_0\to M_n(K_0)\)。对 \(a\in A\)，暂记
\[
p_a(T)=\det\bigl(TI_n-\rho(a\otimes1)\bigr).
\]
这是 \(K_0[T]\) 中的首一 \(n\) 次多项式。

\begin{theorem}\label{b2:thm:reduced-characteristic-descent}
设 \(k\) 为域，\(A\) 为次数 \(n\) 的中心单 \(k\) 代数，\(a\in A\)。对任意分裂域 \(K/k\) 及 \(K\) 代数同构 \(\rho:A\otimes_kK\rightarrow M_n(K)\)，令
\[
p_a(T)=\det\bigl(TI_n-\rho(a\otimes1)\bigr).
\]
则 \(p_a(T)\) 的系数属于 \(k\)，并且与分裂域 \(K\) 及同构 \(\rho\) 的选择无关。
\end{theorem}
\begin{proof}
对 \(g\in\operatorname{Gal}(K_0/k)\)，记 \(g_M\) 为对矩阵条目逐项施加 \(g\) 的映射。直接复合 \(g_M\circ\rho\) 得到的是 \(g\)-半线性映射，不能与 \(\rho\) 作为 \(K_0\) 代数同构比较。为抵消这一半线性，定义
\[
\rho_g=g_M\circ\rho\circ(\operatorname{id}_A\otimes g^{-1}):
A\otimes_kK_0\longrightarrow M_n(K_0).
\]
首尾两个映射分别为 \(g\)-半线性及 \(g^{-1}\)-半线性，复合后为 \(K_0\) 线性；三个映射均保持乘法且双射，所以 \(\rho_g\) 是 \(K_0\) 代数同构。特别地，对任意有限和有
\[
\rho_g\left(\sum_i a_i\otimes\lambda_i\right)
=\sum_i\lambda_i\,g_M\bigl(\rho(a_i\otimes1)\bigr).
\]
于是 \(\rho_g\rho^{-1}\) 是全矩阵代数的 \(K_0\) 自同构，由\cref{b2:thm:skolem-noether}为内自同构。记 \(B=\rho(a\otimes1)\)，则存在 \(U_g\in\operatorname{GL}_n(K_0)\) 使 \(g_M(B)=U_gBU_g^{-1}\)。令 \(g\) 固定不定元 \(T\)，便有
\[
g\bigl(p_a(T)\bigr)
=\det\bigl(TI_n-g_M(B)\bigr)
=\det\bigl(TI_n-U_gBU_g^{-1}\bigr)
=p_a(T).
\]
因此全部系数属于不动域 \(K_0^{\operatorname{Gal}(K_0/k)}=k\)。

固定 \(K_0\) 而改变 \(\rho\)，同样由内自同构的结论，矩阵只差共轭，故多项式不变。再给定任意分裂域 \(L/k\) 及其分裂同构。取 \(L\) 的代数闭包 \(\Omega=\overline L\)。由于 \(K_0/k\) 是有限代数扩张，可逐次为其生成元在 \(\Omega\) 中选择最小多项式的根，将 \(k\hookrightarrow\Omega\) 延拓为 \(K_0\hookrightarrow\Omega\)。这样就得到同时容纳 \(K_0\) 与 \(L\) 的域。将两个分裂模型都扩张到 \(\Omega\)，它们成为 \(A_\Omega\) 到 \(M_n(\Omega)\) 的两个同构，所以仍只差内自同构。因此两个特征多项式在 \(\Omega[T]\) 中相同。第一模型的系数已属于 \(k\)，于是第二模型也给出同一多项式。
\end{proof}

\begin{definition}\label{b2:def:reduced-trace-norm}
设 \(k\) 为域，\(A\) 为次数 \(n\) 的中心单 \(k\) 代数，\(a\in A\)。取 \(A\) 的任意分裂域 \(K/k\) 和 \(K\) 代数同构 \(\rho:A\otimes_kK\rightarrow M_n(K)\)。由\cref{b2:thm:reduced-characteristic-descent}，多项式 \(\det\bigl(TI_n-\rho(a\otimes1)\bigr)\in k[T]\) 与这些选择无关；称它为 \(a\) 的\term[yuehua-tezheng-duoxiangshi]{约化特征多项式}[reduced characteristic polynomial]，记为 \(\chi^{\mathrm{red}}_{A,a}(T)\)。写成
\[
\chi^{\mathrm{red}}_{A,a}(T)=T^n-c_1(a)T^{n-1}+\cdots+(-1)^n c_n(a),
\]
定义\term[yuehua-ji]{约化迹}[reduced trace]与\term[yuehua-fanshu]{约化范数}[reduced norm]为
\[
\operatorname{Trd}_{A/k}(a)=c_1(a),\qquad
\operatorname{Nrd}_{A/k}(a)=c_n(a).
\]
\end{definition}
\symidx{\(\chi^{\mathrm{red}}_{A,a}\)：约化特征多项式}
\symidx{\(\operatorname{Trd}_{A/k}\)：约化迹}
\symidx{\(\operatorname{Nrd}_{A/k}\)：约化范数}

分裂情形 \(A=M_n(k)\) 中，这就是通常的特征多项式、矩阵迹与行列式。约化特征多项式的系数属于中心 \(k\)，所以可以在非交换代数 \(A\) 中将 \(T\) 代为 \(a\)，不必为系数与 \(a\) 的乘法另作次序选择。\cref{b2:thm:exterior-cayley-hamilton}在分裂模型中给出 \(\chi^{\mathrm{red}}_{A,a}(a)=0\)：其像为零，而 \(A\to A_K\) 单射，故原代数中的值也为零。

\subsection{与域迹、域范数的比较}
\label{b2:subsec:1803:02}

令 \(\lambda_a:A\to A\) 表示左乘 \(a\)，将它看成 \(n^2\) 维 \(k\) 空间上的线性算子。其普通迹和行列式与约化运算并不相同。

\begin{proposition}\label{b2:prop:regular-reduced-characteristic}
设 \(k\) 为域，\(A\) 为次数 \(n\) 的中心单 \(k\) 代数，\(a\in A\)。令 \(\lambda_a:A\rightarrow A\) 为左乘 \(a\) 的 \(k\) 线性算子，则其特征多项式满足
\[
\chi_{\lambda_a}(T)=\bigl(\chi^{\mathrm{red}}_{A,a}(T)\bigr)^n.
\]
因此 \(\operatorname{tr}_k\lambda_a=n\operatorname{Trd}_{A/k}(a)\)，且 \(\det_k\lambda_a=\operatorname{Nrd}_{A/k}(a)^n\)。
\end{proposition}
\begin{proof}
扩张到分裂域后，左乘矩阵 \(B=\rho(a\otimes1)\) 逐列作用在 \(M_n(K)\) 上：每一列都是一份 \(K^n\)，都受到同一个 \(B\) 的作用，列与列之间互不混合。因此这个 \(n^2\) 维线性作用是 \(n\) 份 \(B\) 的直和，特征多项式为 \(\det(TI-B)^n\)。标量扩张不改变一个原系数矩阵的特征多项式，因而等式在 \(k[T]\) 中成立。比较次高次系数和常数项，就得到两个公式。
\end{proof}

若底域特征整除 \(n\)，不能把正则迹除以 \(n\) 来定义约化迹。例如在 \(M_p(k)\)、\(\operatorname{char}k=p\) 中，\(E_{11}\) 的约化迹为一，而它的左乘算子迹为零。定义约化运算所需的下降论证正好避开了这种除法。

\begin{theorem}\label{b2:thm:reduced-field-comparison}
设 \(k\) 为域，\(A\) 为次数 \(n\) 的中心单 \(k\) 代数，\(E\subseteq A\) 为含 \(k\) 的子域，\([E:k]=e\)，并按左乘将 \(A\) 视为左 \(E\) 向量空间。则 \(e\mid n\)。对 \(a\in E\)，令 \(m_a:E\rightarrow E\) 为乘 \(a\) 的 \(k\) 线性算子，则有
\[
\chi^{\mathrm{red}}_{A,a}(T)=\chi_{m_a:E\to E}(T)^{n/e},
\]
并且
\[
\operatorname{Trd}_{A/k}(a)=\frac ne\operatorname{Tr}_{E/k}(a),\qquad
\operatorname{Nrd}_{A/k}(a)=\operatorname{N}_{E/k}(a)^{n/e}.
\]
此处 \(E/k\) 不必可分。特别地，次数为 \(n\) 的子域上的约化迹、范数恰为域迹、范数。
\end{theorem}
\begin{proof}
\cref{b2:prop:csa-self-centralizing-field}已证明任意含 \(k\) 的子域 \(E\subseteq A\) 的次数 \(e\) 整除 \(n\)，这里不要求 \(E/k\) 可分。

\(A\) 作为左 \(E\) 空间的维数为 \(n^2/e\)。选一组左 \(E\) 基 \(v_1,\ldots,v_{n^2/e}\)，则左乘 \(a\in E\) 把 \(\sum_j c_jv_j\) 送到 \(\sum_j (ac_j)v_j\)。每个系数空间都是一份 \(E\)，其底层 \(k\) 作用为 \(m_a\)，故同一个算子 \(\lambda_a\) 的特征多项式为 \(\chi_{m_a:E\to E}^{n^2/e}\)。另一方面，\cref{b2:prop:regular-reduced-characteristic}通过分裂矩阵的列坐标给出它的特征多项式为 \((\chi^{\mathrm{red}}_{A,a})^n\)。比较这两种坐标计算，得到
\[
\left(\chi^{\mathrm{red}}_{A,a}\right)^n
=\left(\chi_{m_a:E\to E}^{n/e}\right)^n.
\]
两个括号中的多项式都首一。在 \(k[T]\) 的唯一分解中，每个首一不可约因子在两边的指数分别为 \(n u\) 和 \(n v\)，所以 \(u=v\)，从而两个括号相等。这是在整数中比较因子指数，并未在 \(k\) 中除以 \(n\)。域迹和域范数按乘法算子的迹与行列式定义，见第一卷定义18.2.1，故比较系数得到所述公式。
\end{proof}

非可分子域上的限制可能丢失迹配对的信息。例如取 \(k=\mathbb F_p(t)\) 和 \(E=k(\alpha)\)，其中 \(p\) 为素数，\(\alpha^p=t\)。多项式 \(T^p-t\in\mathbb F_p[t][T]\) 对素元 \(t\) 满足第一卷定理12.3.1的 Eisenstein 判据，所以它在 \(k[T]\) 中不可约，\([E:k]=p\)，且 \(1,\alpha,\ldots,\alpha^{p-1}\) 为一组基。在这组基下，正则作用 \(a\mapsto m_a\) 将 \(E\) 嵌入 \(A=\operatorname{End}_k(E)\cong M_p(k)\)。对 \(1\le i<p\)，乘 \(\alpha^i\) 将每个基向量送到另一个基向量的标量倍数，没有对角项；乘标量 \(c\) 的迹则为 \(pc=0\)。由迹的线性，\(\operatorname{Tr}_{E/k}(a)=0\) 对所有 \(a\in E\) 成立。比较定理中的 \(e=n=p\)，所以 \(\operatorname{Trd}_{A/k}\) 在 \(E\) 上恒为零。由于 \(xy\in E\) 对 \(x,y\in E\) 成立，配对 \((x,y)\mapsto\operatorname{Trd}_{A/k}(xy)\) 在 \(E\times E\) 上也恒为零；\cref{b2:prop:involution-reduced-trace-form}将证明它在全代数 \(A\) 上非退化。

\ProseParagraphBreak
在特征不为二的域 \(k\) 上，取 \(a,b\in k^\times\) 及四元数代数 \(Q=(a,b)_k\)。\cref{b2:prop:quaternion-norm}中的共轭与范数给出每个 \(x\in Q\) 的关系
\[
x^2-(x+\bar x)x+x\bar x=0.
\]
其中 \(x+\bar x\) 与 \(x\bar x\) 都属于 \(k\)。若 \(x\) 非标量，它不能满足非零的一次 \(k\) 系数关系，否则就属于 \(k\)；所以最小多项式 \(\mu_x(T)\) 恰为这个首一二次多项式。约化 Cayley--Hamilton 等式又使 \(\mu_x\) 整除首一二次多项式 \(\chi^{\mathrm{red}}_{Q,x}\)，故两者相等。若 \(x=c\in k\)，则 \(\mu_c(T)=T-c\)，而分裂矩阵为 \(cI_2\)，所以 \(\chi^{\mathrm{red}}_{Q,c}(T)=(T-c)^2\)。由这两种情形比较系数，约化迹与范数分别为 \(x+\bar x\) 与 \(x\bar x\)，标量情形即为 \(2c\) 与 \(c^2\)。

\subsection{定义的良定性与基本运算性质}
\label{b2:subsec:1803:03}

\begin{proposition}\label{b2:prop:reduced-trace-norm-properties}
设 \(k\) 为域，\(A\) 为次数 \(n\) 的中心单 \(k\) 代数，则约化迹 \(\operatorname{Trd}_{A/k}:A\rightarrow k\) 为 \(k\) 线性映射；对任意 \(a,b\in A\)，约化范数 \(\operatorname{Nrd}_{A/k}:A\rightarrow k\) 满足
\[
\operatorname{Nrd}(ab)=\operatorname{Nrd}(a)\operatorname{Nrd}(b),\qquad
\operatorname{Nrd}(1)=1.
\]
对 \(c\in k\)，有 \(\operatorname{Trd}(c1_A)=nc\)、\(\operatorname{Nrd}(c1_A)=c^n\)。元素 \(a\) 可逆，当且仅当 \(\operatorname{Nrd}(a)\ne0\)。
\end{proposition}
\begin{proof}
在一个分裂域中计算，迹的线性、行列式的乘法性与标量矩阵公式给出前面各项；值已由\cref{b2:thm:reduced-characteristic-descent}属于 \(k\)，故等式就在底域中成立。

若 \(a\) 可逆，乘法性给出 \(\operatorname{Nrd}(a)\operatorname{Nrd}(a^{-1})=1\)。反过来，范数非零使分裂矩阵可逆，故逐列作用的 \(\lambda_a\) 扩域后可逆。它的原 \(k\) 矩阵的行列式在扩域中非零，因而在 \(k\) 中也非零，这就把矩阵的可逆性降回了底域。于是存在 \(b\in A\) 使 \(ab=1\)；再由 \(a(ba-1)=0\) 及 \(\lambda_a\) 单射得 \(ba=1\)。因此 \(a\) 是单位。
\end{proof}

约化范数的乘法性不包含可加性。取 \(A=M_2(k)\)，则 \(E_{11}\) 与 \(E_{22}\) 的行列式均为零，而 \(E_{11}+E_{22}=I_2\) 的行列式为一，所以 \(\operatorname{Nrd}(E_{11}+E_{22})\ne\operatorname{Nrd}(E_{11})+\operatorname{Nrd}(E_{22})\)。

\ProseParagraphBreak
对原元素 \(a\otimes1\)，扩域后的约化系数可以在共同的分裂模型中比较。但 \(A_F\) 还包含坐标不属于 \(k\) 的新元素 \(\sum_i a_i\otimes x_i\)。要同时处理这些元素，需要证明约化系数本身是坐标的 \(k\) 系数多项式，而不是只比较原元素的取值。

\begin{proposition}\label{b2:prop:reduced-base-change}
设 \(k\) 为域，\(A\) 为次数 \(n\) 的中心单 \(k\) 代数，\(F/k\) 为域扩张，则约化特征多项式及其各系数与标量扩张 \(A_F=A\otimes_kF\) 相容。更具体地，取 \(A\) 的 \(k\) 基 \(a_1,\ldots,a_{n^2}\)，则元素 \(\sum_i x_i a_i\in A\) 的约化特征多项式的每个系数都是坐标 \(x_1,\ldots,x_{n^2}\) 的 \(k\) 系数多项式；对 \(A_F\) 中的 \(\sum_i a_i\otimes x_i\)、\(x_i\in F\)，仍由同一个多项式计算。
\end{proposition}
\begin{proof}
选有限 Galois 分裂域 \(K/k\) 及分裂同构 \(\rho\)，令
\[
H(\mathbf X)=\sum_i X_i\rho(a_i\otimes1).
\]
其特征多项式的系数属于 \(K[X_1,\ldots,X_{n^2}]\)。在\cref{b2:thm:reduced-characteristic-descent}的证明中，\(\rho_g\rho^{-1}\) 是整个矩阵代数的一个内自同构，所以比较矩阵 \(U_g\) 与输入元素无关。令 \(g\) 固定所有不定元 \(X_i\)，便有
\[
g\bigl(H(\mathbf X)\bigr)=U_gH(\mathbf X)U_g^{-1}.
\]
因此每个特征系数作为多项式都被 \(\operatorname{Gal}(K/k)\) 固定。逐项比较单项式的标量系数，得到这些系数属于 \(k\)，即特征系数多项式属于 \(k[X_1,\ldots,X_{n^2}]\)。

给定 \(F/k\)，如下降定理的证明那样，将有限代数扩张 \(K/k\) 嵌入 \(F\) 的代数闭包 \(\Omega\)。上述分裂同构扩张后，也是 \((A_F)_\Omega\cong A_\Omega\) 的分裂模型。在此模型中将 \(X_i\) 代为任意 \(x_i\in F\)，就得到 \(\sum_i a_i\otimes x_i\) 的特征多项式。所有系数都由刚得到的 \(k\) 系数多项式取值，故属于 \(F\)；分裂模型选择无关性保证它们就是 \(A_F\) 的约化系数。
\end{proof}

特别地，对原来的 \(a\in A\)，在任何扩张中都有
\[
\operatorname{Trd}_{A_F/F}(a\otimes1)=\operatorname{Trd}_{A/k}(a),\qquad
\operatorname{Nrd}_{A_F/F}(a\otimes1)=\operatorname{Nrd}_{A/k}(a).
\]
泛型矩阵的行列式每一项都恰含 \(n\) 个矩阵条目，所以约化范数的坐标多项式是 \(n\) 次齐次多项式。

\begin{proposition}\label{b2:prop:involution-reduced-trace-form}
设 \(k\) 为域，\(A\) 为次数 \(n\) 的中心单 \(k\) 代数。则
\[
\beta(x,y)\coloneqq\operatorname{Trd}_{A/k}(xy)
\]
是 \(A\) 上非退化的对称 \(k\) 双线性型，且 \(\operatorname{Trd}_{A/k}:A\to k\) 是满射。

若 \(\sigma:A\to A\) 为一个 \(k\) 线性映射，满足 \(\sigma(1)=1\)、\(\sigma^2=\operatorname{id}\)，且对所有 \(x,y\in A\) 有 \(\sigma(xy)=\sigma(y)\sigma(x)\)，则 \(\sigma(x)\) 与 \(x\) 有相同的约化特征多项式，因而约化迹与约化范数均相同。并且
\[
t_\sigma(x,y)\coloneqq\operatorname{Trd}_{A/k}\bigl(\sigma(x)y\bigr)
\]
是 \(A\) 上非退化的对称 \(k\) 双线性型，称为 \(\sigma\) 的\term[duihe-ji-xing]{迹型}[trace form of an involution]。
\end{proposition}
\begin{proof}
约化迹的线性给出 \(\beta\) 的双线性。选分裂域 \(K/k\) 及分裂同构 \(\rho\)，则 \(\beta\) 在分裂模型中由 \((X,Y)\mapsto\operatorname{tr}(XY)\) 给出。矩阵迹满足 \(\operatorname{tr}(XY)=\operatorname{tr}(YX)\)，故 \(\beta(x,y)=\beta(y,x)\)。这个矩阵配对还非退化：若 \(X_{ij}\ne0\)，取 \(Y=E_{ji}\)，便有 \(\operatorname{tr}(XY)=X_{ij}\ne0\)。

取 \(A\) 的一组 \(k\) 基 \(a_1,\ldots,a_{n^2}\)，其像 \(\rho(a_i\otimes1)\) 为 \(M_n(K)\) 的 \(K\) 基。原配对的 Gram 矩阵为 \(G=(\operatorname{Trd}(a_i a_j))_{i,j}\in M_{n^2}(k)\)，它在 \(K\) 中就是上述矩阵配对的 Gram 矩阵，因此 \(\det G\ne0\)。这个行列式原属 \(k\)，故 \(\beta\) 在 \(k\) 上也非退化。由于 \(1_A\ne0\)，存在 \(y\in A\) 使 \(\beta(1_A,y)=\operatorname{Trd}(y)\ne0\)。约化迹是到一维空间 \(k\) 的非零线性映射，因而满射。这里不需要 \(\operatorname{Trd}(1_A)=n\) 非零。

若存在所述 \(\sigma\)，将 \(\sigma\otimes\operatorname{id}_K\) 通过 \(\rho\) 转为 \(M_n(K)\) 上的反自同构 \(\tau\)。转置 \(t:X\mapsto X^{\mathsf T}\) 也反转乘法次序，所以 \(\tau\circ t\) 是 \(K\) 代数自同构。由\cref{b2:thm:skolem-noether}，存在 \(P\in\operatorname{GL}_n(K)\) 使 \(\tau(X)=PX^{\mathsf T}P^{-1}\)。因此 \(\sigma(x)\) 的分裂矩阵与 \(x\) 的分裂矩阵之转置相似，特征多项式不变。\cref{b2:thm:reduced-characteristic-descent}把这个等式降回 \(k\)，比较系数就得到约化迹和约化范数不变。

迹型的双线性来自 \(\sigma\) 和约化迹的线性。再把刚证的迹不变性用于 \(\sigma(x)y\)，得到
\[
\operatorname{Trd}\bigl(\sigma(x)y\bigr)
=\operatorname{Trd}\Bigl(\sigma\bigl(\sigma(x)y\bigr)\Bigr)
=\operatorname{Trd}\bigl(\sigma(y)x\bigr),
\]
所以它对称。又因为 \(t_\sigma(x,y)=\beta(\sigma(x),y)\)，而 \(\sigma^2=\operatorname{id}\) 保证 \(\sigma\) 是可逆线性映射，将 \(\beta\) 的第一变量与它复合仍然非退化。
\end{proof}

对实矩阵代数的转置对合，\(t_\sigma(X,X)=\operatorname{tr}(X^{\mathsf T}X)\) 为所有矩阵条目的平方和；对四元数的标准共轭 \(\gamma\)，次数为二的标量迹公式则给出
\[
t_\gamma(x,x)=\operatorname{Trd}(\bar x x)=2N(x).
\]
这将乘法范数与一个双线性型联系起来。迹型依赖所选对合，不只是底层代数的属性；\cref{b2:ex:18-two-involution-trace-forms}会在同一个矩阵代数上比较两种不同的情形。

% !TeX root = ../../Book2.tex
% 纲要标题：### 18.4 交叉积与相对 Brauer 群
% 纲要位置：计划纲要.md，第 1652 行。

\section{交叉积与相对 Brauer 群}
\label{b2:sec:1804}

循环代数只需一个生成自同构和一个参数。一般有限 Galois 扩张有多个自同构；记录它们对应元素相乘时的标量误差，就得到\(2\)-余圈。本节用显式方程定义这些数据，证明它们与相对 Brauer 群的对应；第四卷再从一般上链复形研究群上同调。

% 纲要要点（供目录核对）：
%
% * 交叉积与2-余圈；结合律产生余圈方程，重选基元产生余边；正式证明因子集不依赖标量提升及矩阵模型
% * 有限 Galois 扩张的相对 Brauer 群及循环情形的范数商；对角幂等元对齐两份子域作用，取角证明群运算，同类同次数及 Skolem–Noether 证明单射，明确下降因子集的逆类约定
% * 与第四卷群上同调的衔接；循环余圈化为进位形式，正文证明交叉积约化迹只由单位指标分量贡献
% * Brauer 群不能与第五卷 Brauer 诱导定理混同
%

\subsection{\texorpdfstring{交叉积与 \(2\)-余圈}{交叉积与 2-余圈}}
\label{b2:subsec:1804:01}

固定有限 Galois 扩张 \(K/k\)，记 \(G=\operatorname{Gal}(K/k)\)、\(n=|G|=[K:k]\)。希望为每个 \(g\in G\) 配一个基元 \(u_g\)，使它对标量的作用由 \(g\) 给出，而两个基元相乘只与 \(u_{gh}\) 相差一个非零标量，即
\[
u_g a=g(a)u_g,\qquad u_g u_h=c(g,h)u_{gh}.
\]
这两个规则把结合律转成一个可计算的条件：
\[
(u_g u_h)u_t=c(g,h)c(gh,t)u_{ght},\qquad
u_g(u_h u_t)=g\bigl(c(h,t)\bigr)c(g,ht)u_{ght}.
\]
若再要求 \(u_1=1\)，就得到下面的余圈方程及归一化条件。

\begin{definition}\label{b2:def:normalized-two-cocycle}
设 \(K/k\) 为有限 Galois 扩张，\(G=\operatorname{Gal}(K/k)\)，\(G\) 按域自同构作用于 \(K^\times\)。若函数 \(c:G\times G\to K^\times\) 对任意 \(g,h,t\in G\) 满足
\begin{equation}\label{b2:eq:two-cocycle}
c(g,h)c(gh,t)=g\bigl(c(h,t)\bigr)c(g,ht),\qquad
c(1,g)=c(g,1)=1,
\end{equation}
则称 \(c\) 为\term[guiyihua-er-yuquan]{归一化 \(2\)-余圈}[normalized 2-cocycle]。若取元素族 \((b_g)_{g\in G}\subseteq K^\times\) 且 \(b_1=1\)，定义函数
\[
(\delta b)(g,h)=b_g\cdot g(b_h)b_{gh}^{-1}
\]
并称 \(\delta b\) 为\term[er-yubian]{\(2\)-余边}[2-coboundary]。
\end{definition}

直接将 \(\delta b\) 代入\eqref{b2:eq:two-cocycle}，两边都化为 \(b_g\cdot g(b_h)gh(b_t)b_{ght}^{-1}\)，所以余边确为余圈。

\ProseParagraphBreak
余边记录的是同一组基元的缩放。若将 \(u_g\) 换成 \(u'_g=b_g u_g\)，其中 \(b_1=1\)，则
\[
u'_g u'_h
=b_g g(b_h)c(g,h)b_{gh}^{-1}u'_{gh}
=\bigl(c\delta b\bigr)(g,h)u'_{gh}.
\]
因此改变这些基元会改变乘法系数，但只使余圈乘上一个余边。

\begin{definition}\label{b2:def:galois-crossed-product}
设 \(K/k\) 为有限 Galois 扩张，\(G=\operatorname{Gal}(K/k)\)，\(c:G\times G\rightarrow K^\times\) 为归一化 \(2\)-余圈。在左 \(K\) 空间 \(B_c=\bigoplus_{g\in G}Ku_g\) 上规定
\begin{equation}\label{b2:eq:crossed-product-multiplication}
(a u_g)(b u_h)=a\cdot g(b)c(g,h)u_{gh}.
\end{equation}
按双可加性延拓后，所得含幺 \(k\) 代数称为由 \(K/k\) 和 \(c\) 确定的\term[jiaocha-ji-daishu]{交叉积代数}[crossed product algebra]。
\end{definition}
\symidx{\(B_c=(K/k,G,c)\)：Galois 扩张的交叉积代数}

其中 \(K\) 通过 \(a\mapsto au_1\) 嵌入；乘法公式同时规定了 \(u_g b=g(b)u_g\) 与 \(u_g u_h=c(g,h)u_{gh}\)。前一条控制基元经过标量时发生的自同构，后一条控制基元之间的乘法误差。前述换基计算也给出一个同构 \(B_{c\delta b}\to B_c\)：把新代数中指标 \(g\) 的基元送到 \(b_g u_g\) 即可。

\begin{proposition}\label{b2:prop:crossed-product-central-simple}
设 \(K/k\) 为有限 Galois 扩张，\(G=\operatorname{Gal}(K/k)\)，\(n=[K:k]\)，\(c:G\times G\rightarrow K^\times\) 为归一化 \(2\)-余圈，则相应交叉积 \(B_c=\bigoplus_{g\in G}Ku_g\) 是次数 \(n\) 的中心单 \(k\) 代数，单位为 \(u_1\)。子域 \(K\) 分裂它；每个 \(u_g\) 都可逆。
\end{proposition}
\begin{proof}
比较三项乘积 \((a u_g)(b u_h)(d u_t)\)：除共同因子 \(a\cdot g(b)gh(d)\) 外，两种结合方式的系数分别是\eqref{b2:eq:two-cocycle}的两边。因此乘法结合。归一化条件使 \(u_1\) 为单位，底域标量在中心。向量空间维数为 \(n[K:k]=n^2\)。由余圈条件还有 \(c(g,g^{-1})=g\bigl(c(g^{-1},g)\bigr)\)，故
\[
u_g^{-1}=c(g^{-1},g)^{-1}u_{g^{-1}}.
\]
左右相乘均为一。

中心的计算与循环代数同样，先检查与 \(K\) 交换的条件，再检查与全部基元交换的条件。若 \(x=\sum_g a_g u_g\) 与全部 \(b\in K\) 交换，则比较系数得 \(a_g\bigl(g(b)-b\bigr)=0\)。对 \(g\ne1\) 可选使差非零的 \(b\)，所以只有单位指标的系数可能非零。再与各 \(u_g\) 交换，就使该系数被全部 \(G\) 固定，因而属于 \(k\)。故中心是 \(k\)。

单性的证明沿用循环代数中减少非零项数的方法。对一个非零双边理想，选非零项数最少的元素，并取其中一个非零项的指标 \(g\)。右乘 \(u_g^{-1}\) 使每个指标 \(h\) 变为 \(hg^{-1}\)，各系数只乘上非零标量，所以非零项数不变，而单位指标的系数变为非零。再左乘一个 \(K^\times\) 中的标量，使该系数为一。所得元素与任意 \(b\in K\) 的交换子在单位指标处为零，其他非零项只可能来自原有的指标；若交换子非零，就有更少的非零项，违反最少性。于是该元素与 \(K\) 交换，由刚才的系数比较只能为一，理想即为全代数。这证明单性。最后 \([K:k]=n=\deg B_c\)，由\cref{b2:prop:maximal-field-splits-csa}，\(K\) 分裂 \(B_c\)。
\end{proof}

若交叉积分裂为 \(n\) 阶矩阵代数，它在一个 \(n\) 维 \(k\) 空间上的作用限制到 \(K\)，就成为一维 \(K\) 空间上的作用。在这个空间上，每个 \(u_g\) 的作用满足 \(T_g(az)=g(a)T_g(z)\)；选择一个 \(K\) 基以后，它只能是域自同构 \(g\) 再乘上一个非零标量。这些标量正是余边公式中的 \(b_g\)。

\begin{proposition}\label{b2:prop:crossed-product-split-coboundary}
设 \(K/k\) 为有限 Galois 扩张，\(G=\operatorname{Gal}(K/k)\)，\(c:G\times G\rightarrow K^\times\) 为归一化 \(2\)-余圈，\(B_c\) 为相应交叉积代数，则 \(B_c\) 分裂，当且仅当 \(c\) 是 \(2\)-余边。
\end{proposition}
\begin{proof}
若 \(c(g,h)=b_g\cdot g(b_h)b_{gh}^{-1}\)，其中 \(b_1=1\)，在 \(k\) 空间 \(K\) 上让 \(K\) 左乘，让 \(u_g\) 按 \(T_g(z)=b_g\cdot g(z)\) 作用。计算得到
\[
T_gT_h(z)=b_g g(b_h)gh(z)=c(g,h)T_{gh}(z),\qquad
T_g(az)=g(a)T_g(z).
\]
因此这些算子满足交叉积的两个乘法规则，且 \(T_1=\operatorname{id}_K\)，给出含幺同态 \(B_c\to\operatorname{End}_k(K)\)。由单性及两端同为 \(n^2\) 维，它是同构，其中 \(n=[K:k]\)。

若 \(B_c\cong\operatorname{End}_k(W)\)，由 \(\deg B_c=n\) 有 \(\dim_k W=n\)。限制到 \(K\) 后，维数公式给出 \(\dim_KW=1\)。选一个 \(K\) 基，将 \(W\) 识别为 \(K\)。由于 \(u_g z=g(z)u_g\)，其作用满足 \(T_g(z)=g(z)T_g(1)=b_g g(z)\)，其中 \(b_g=T_g(1)\ne0\)，且 \(b_1=1\)。代入 \(u_g u_h=c(g,h)u_{gh}\)，得 \(b_g\cdot g(b_h)=c(g,h)b_{gh}\)，正是余边条件。此时把 \(u_g\) 换成 \(b_g^{-1}u_g\)，它们在 \(K\) 上就按 \(g\) 作用，乘法误差也同时消失。
\end{proof}

\subsection{Brauer 群与群上同调}
\label{b2:subsec:1804:02}

\(2\)-余圈按逐点相乘组成交换群，余边构成其子群。定义商群
\[
H^2(G,K^\times)=\{\text{归一化 }2\text{-余圈}\}/\{2\text{-余边}\}.
\]
这个商群是当前 \(G\) 作用下二阶\term[qun-shangtongdiao]{群上同调}[group cohomology]的显式模型。第四卷第 12.1 节从 bar 消解建立余链复形，第 12.2 节再用低阶余圈解释群扩张；这里把交换群 \(K^\times\) 的运算写成乘法，得到同一套二阶方程，下面的论证只用这些显式方程。
\symidx{\(H^2(G,K^\times)\)：乘法系数的二阶群上同调}

\ProseParagraphBreak
对于 \(g\in G\)，可加映射 \(T_g\) 若满足 \(T_g(av)=g(a)T_g(v)\)，称其为 \(g\)-\term[banxianxing-yingshe]{半线性映射}[semilinear map]。在矩阵代数上，逐条目施加 \(g\) 就给出这种映射。下面把一个已分裂代数的下降数据写成具体矩阵。

\begin{proposition}\label{b2:prop:split-algebra-factor-set}
设 \(K/k\) 为有限 Galois 扩张，\(G=\operatorname{Gal}(K/k)\)，\(A\) 为由 \(K\) 分裂的中心单 \(k\) 代数，\(m=\deg A\)，取 \(K\) 代数同构 \(\rho:A_K=A\otimes_kK\rightarrow M_m(K)\)。令 \(\alpha_g=\rho(\operatorname{id}_A\otimes g)\rho^{-1}\) 为 \(A_K\) 的自然 Galois 作用在矩阵代数上的表示，\(g(X)\) 表示对矩阵 \(X\) 逐条目施加 \(g\)，则存在 \(C_g\in\operatorname{GL}_m(K)\)、\(C_1=I_m\)，使 \(\alpha_g(X)=C_g\cdot g(X)C_g^{-1}\)。由
\[
C_g\cdot g(C_h)=c(g,h)C_{gh}
\]
定义的标量 \(c(g,h)\in K^\times\) 构成归一化 \(2\)-余圈，而且 \([A]=-[B_c]\)，其中 \(B_c\) 是相应交叉积代数。若将 \(C_g\) 换为 \(b_gC_g\)，其中 \(b_g\in K^\times\)、\(b_1=1\)，则余圈换为 \(c\delta b\)。若将分裂模型换为 \(\rho'=\operatorname{Ad}(P)\rho\)，其中 \(P\in\operatorname{GL}_m(K)\)、\(\operatorname{Ad}(P)(X)=PXP^{-1}\)，并取
\[
C'_g=PC_g g(P)^{-1},
\]
则得到原余圈 \(c\)。特别地，所得类 \([c]\in H^2(G,K^\times)\) 与 \(C_g\) 及分裂同构 \(\rho\) 的选择无关。
\end{proposition}
\begin{proof}
\(\alpha_g\) 在标量上作用为 \(g\)，即为 \(g\)-半线性代数自同构。将它与逐条目作用 \(g^{-1}\) 依次复合，所得 \(\alpha_g\circ g^{-1}\) 是一个 \(K\) 代数自同构；由\cref{b2:thm:skolem-noether}，存在 \(C_g\in\operatorname{GL}_m(K)\)，使
\begin{equation}\label{b2:eq:brauer-semilinear-matrices}
\alpha_g(X)=C_g\cdot g(X)C_g^{-1},\qquad C_1=I_m.
\end{equation}
\(\alpha_g\alpha_h=\alpha_{gh}\) 说明 \(C_g\cdot g(C_h)C_{gh}^{-1}\) 与所有矩阵交换，所以为唯一的标量 \(c(g,h)\in K^\times\)。比较同一个矩阵积的两种结合方式，得到
\[
\begin{aligned}
\bigl(C_g g(C_h)\bigr)gh(C_t)
&=c(g,h)c(gh,t)C_{ght},\\
C_g g\bigl(C_h h(C_t)\bigr)
&=g\bigl(c(h,t)\bigr)c(g,ht)C_{ght}.
\end{aligned}
\]
消去可逆矩阵 \(C_{ght}\)，即为\eqref{b2:eq:two-cocycle}；\(C_1=I_m\) 又给出归一化条件。

对同一个 \(\alpha_g\)，两个实现矩阵只相差非零标量，因为它们的商与全部矩阵交换。因此任意另一组归一化的实现矩阵都形如 \(b_gC_g\)、\(b_1=1\)，而
\[
(b_gC_g)g(b_hC_h)
=b_g g(b_h)c(g,h)C_{gh}
=\bigl(c\delta b\bigr)(g,h)(b_{gh}C_{gh}).
\]
再改变分裂模型。由 \(\rho'=\operatorname{Ad}(P)\rho\) 得
\[
\alpha'_g(X)
=P C_g g(P)^{-1}g(X)g(P)C_g^{-1}P^{-1}.
\]
所述 \(C'_g\) 因而实现 \(\alpha'_g\)，且 \(C'_1=I_m\)。它们满足
\[
C'_g g(C'_h)
=PC_g g(C_h)gh(P)^{-1}
=c(g,h)C'_{gh},
\]
所以余圈不变。任意两个 \(K\) 代数分裂同构之差都是 \(M_m(K)\) 的 \(K\) 自同构，由 Skolem--Noether 为这种内自同构；这就涵盖了全部分裂模型的选择。

在 \(k\) 空间 \(V=K^m\) 上，让 \(K\) 按标量乘法作用，让 \(u_g\) 按 \(v\mapsto C_g\cdot g(v)\) 作用，得到 \(B_c\) 的表示。它含幺且 \(B_c\) 单，因而忠实。与 \(K\) 作用交换的算子恰好是 \(K\) 线性算子，可写为 \(X\in M_m(K)\)；这样的 \(X\) 与所有 \(u_g\) 作用交换，当且仅当
\[
XC_g=C_g g(X),\qquad\text{即 }X=\alpha_g(X)\quad(g\in G).
\]
这些不动矩阵恰为 \(\rho(A\otimes1)\)：在 \(A_K\) 中按 \(A\) 的一组 \(k\) 基展开，自然 Galois 作用的不动性恰好使全部系数属于 \(K^G=k\)。因此 \(B_c\) 在 \(\operatorname{End}_k(V)\) 中的中心化子同构于 \(A\)。由\cref{b2:cor:csa-centralizer-factorization}的中心为 \(k\) 的情形，
\[
B_c\otimes_k A\cong\operatorname{End}_k(V).
\]
所以 \([A]=-[B_c]\)，得到所述关系。
\end{proof}

\begin{theorem}\label{b2:thm:relative-brauer-cocycles}
设 \(K/k\) 为有限 Galois 扩张，\(G=\operatorname{Gal}(K/k)\)。记 \(H^2(G,K^\times)\) 为归一化 \(2\)-余圈按\(2\)-余边取的商群，\(\operatorname{Br}(K/k)=\ker\bigl(\operatorname{Br}(k)\rightarrow\operatorname{Br}(K)\bigr)\)，则交叉积给出群同构
\[
H^2(G,K^\times)\longrightarrow\operatorname{Br}(K/k),
\qquad[c]\longmapsto[B_c].
\]
\end{theorem}
\begin{proof}
首先，\cref{b2:prop:crossed-product-central-simple}使像落在相对 Brauer 群；换基公式保证它不依赖余圈代表。

为核对乘法，取两个余圈 \(c,d\)，并记 \(n=[K:k]\)。\(B_c\otimes_kB_d\) 的 \(k\) 维数为 \(n^4\)，而 \(B_{cd}\) 的维数为 \(n^2\)；当 \(n>1\) 时两者尺寸不同，一般不能期待它们直接同构。要比较的是 Brauer 类，需要去掉不影响除代数代表的矩阵大小；下面通过非零幂等元取角来实现这一点。有限 Galois 扩张具有代数同构
\begin{equation}\label{b2:eq:galois-tensor-diagonal}
K\otimes_kK\longrightarrow\prod_{g\in G}K,
\qquad a\otimes b\longmapsto\bigl(a\cdot g(b)\bigr)_{g\in G}.
\end{equation}
具体地，第一卷定理16.4.2给出本原元素 \(K=k(\theta)\)；其可分最小多项式在 \(K\) 中分成互异因子 \(T-g(\theta)\)。对 \(K[T]/(f)\) 应用第一卷定理10.4.3的中国剩余定理，便得到上述同构。

选择单位指标，是因为该坐标把 \(a\otimes b\) 送到 \(ab\)：两份 \(K\) 在这一分量上的作用被对齐。记 \(e_t\in K\otimes_kK\) 为在指标 \(t\) 处坐标为一、其他坐标为零的幂等元，并取 \(e=e_1\)，就有 \(e(a\otimes1)e=e(1\otimes a)e\)。因此希望取角以后只保留对两份 \(K\) 同步施加同一自同构的项。在 \(E=B_c\otimes_kB_d\) 中，用 \(u_g,v_h\) 表示两侧的基元，记 \(z_{g,h}=u_g\otimes v_h\)。共轭 \(z_{g,h}\) 对 \(K\otimes_kK\) 作用为 \(g\otimes h\)。若 \(\lambda_t(a\otimes b)=a\,t(b)\)，则
\[
\lambda_t\bigl((g\otimes h)(a\otimes b)\bigr)
=g(a)\,th(b)
=g\bigl(\lambda_{g^{-1}th}(a\otimes b)\bigr).
\]
按可加性，这个等式对 \(K\otimes_kK\) 的全部元素成立。将它用于 \(e\)，可知 \(z_{g,h}e z_{g,h}^{-1}=e_{gh^{-1}}\)。不同坐标幂等元的乘积为零，因而
\[
e z_{g,h}e=e e_{gh^{-1}}z_{g,h}=0\quad\text{若 }g\ne h.
\]
当 \(g=h\) 时，\(z_{g,g}\) 与 \(e\) 交换，令 \(w_g=e z_{g,g}e=e z_{g,g}\)。取角后，系数环 \(e(K\otimes_kK)e\) 经 \(\lambda_1\) 识别为 \(K\)。原代数是自由左 \(K\otimes_kK\) 模，以全部 \(z_{g,h}\) 为基；取角消去了 \(g\ne h\) 的项，留下
\[
eEe=\bigoplus_{g\in G}e(K\otimes_kK)w_g
\cong\bigoplus_{g\in G}K w_g
\]
的左 \(K\) 模分解。这里各项仍然独立，因为它们在原自由模中有不同的指标；而 \(w_g\ne0\)，否则右乘 \(z_{g,g}^{-1}\) 就会使 \(e=0\)。计算乘法时，\(g\otimes g\) 在单位坐标上正是 \(g\)，并且
\[
w_g w_h=e\bigl(c(g,h)\otimes d(g,h)\bigr)z_{gh,gh}
=c(g,h)d(g,h)w_{gh}.
\]
同样，若用 \(e(b\otimes1)e\) 表示角代数中的标量 \(b\in K\)，便有
\[
w_g b=e(g(b)\otimes1)z_{g,g}=g(b)w_g.
\]
于是把 \(B_{cd}\) 的基元送到 \(w_g\)，把标量送到 \(e(K\otimes_kK)e\) 中的对应元素，得到含幺 \(k\) 代数同构 \(B_{cd}\to eEe\)，角代数的单位是 \(e=w_1\)。由\cref{b2:prop:brauer-division-representative}，中心单代数的非零幂等角与原代数有同一 Brauer 类。因此 \([B_c]+[B_d]=[B_{cd}]\)，映射保持群运算。\(c=1\) 时的交叉积由\cref{b2:prop:crossed-product-split-coboundary}分裂，确实对应零元。

若 \([B_c]=[B_d]\)，由\cref{b2:prop:brauer-division-representative}，可写 \(B_c\cong M_r(D)\)、\(B_d\cong M_s(D)\)，其中 \(D\) 为同一个中心除代数。两代数的次数同为 \(n\)，所以 \(r\deg D=s\deg D=n\)，从而 \(r=s\)，可取 \(k\) 代数同构 \(\phi:B_c\to B_d\)。记 \(\iota_c,\iota_d\) 为两代数中 \(K\) 的标准嵌入。由\cref{b2:thm:skolem-noether}，有 \(z\in B_d^\times\)，使
\[
\iota_d(a)=z\phi\bigl(\iota_c(a)\bigr)z^{-1}\qquad(a\in K).
\]
于是 \(\psi=\operatorname{Ad}(z)\phi\) 在两份标准 \(K\) 之间为恒等。将 \(x=\psi(u_g)\) 展开为 \(\sum_h b_{g,h}v_h\)，由 \(xa=g(a)x\) 得
\[
b_{g,h}h(a)=g(a)b_{g,h}\qquad(a\in K).
\]
若 \(b_{g,h}\ne0\)，可消去这个 \(K\) 中的系数，得到 \(h(a)=g(a)\) 对全部 \(a\in K\) 成立，故 \(h=g\)。因此 \(\psi(u_g)=b_gv_g\)，且因 \(u_g\) 可逆，有 \(b_g\in K^\times\)；单位条件给出 \(b_1=1\)。比较 \(\psi(u_g u_h)\) 与 \(\psi(u_g)\psi(u_h)\) 的系数，得到
\[
c(g,h)b_{gh}=b_g\cdot g(b_h)d(g,h).
\]
因此 \(c/d\) 为余边，证明单射性。

最后，对任意 \([A]\in\operatorname{Br}(K/k)\)，\(A_K\) 的 Brauer 类为零，由\cref{b2:prop:brauer-division-representative}在底域 \(K\) 上的结论，它同构于一个 \(K\) 上的全矩阵代数。因此可应用\cref{b2:prop:split-algebra-factor-set}，得到余圈 \(c\)，使 \([A]=-[B_c]\)。由已经证明的乘法相容性，\([B_{c^{-1}}]=-[B_c]\)，所以原类是 \([c^{-1}]\) 的像。这证明满射性，所需群同构成立。
\end{proof}

半线性映射作为 \(k\) 线性算子参与交叉积的表示，却同时改变 \(K\) 标量。乘法公式中的 \(g\bigl(c(h,t)\bigr)\) 正记录了这种改变。从余圈构造交叉积，与从分裂代数的半线性作用提取因子集，这两种构造的关系为
\[
[c]\longmapsto[B_c],\qquad
A\longmapsto c_A,\qquad
[A]=-[B_{c_A}]=[B_{c_A^{-1}}].
\]
前一个箭头是\cref{b2:thm:relative-brauer-cocycles}的同构；从一个已分裂的代数 \(A\) 出发按\cref{b2:prop:split-algebra-factor-set}构造因子集时，原类对应的是 \([c_A^{-1}]\)。

\begin{corollary}[循环情形的范数商]\label{b2:cor:cyclic-relative-brauer}
设 \(K/k\) 为次数 \(n\ge2\) 的循环 Galois 扩张，\(G=\operatorname{Gal}(K/k)=\langle\sigma\rangle\)，则有群同构
\[
H^2(G,K^\times)\cong\operatorname{Br}(K/k)
\cong k^\times/\operatorname{N}_{K/k}(K^\times).
\]
具体地，映射 \(a\operatorname{N}_{K/k}(K^\times)\mapsto[(K/k,\sigma,a)]\) 给出范数商到相对 Brauer 群的同构。
\end{corollary}
\begin{proof}
由\cref{b2:thm:relative-brauer-cocycles}，每个相对 Brauer 类都由某个交叉积 \(B_c\) 表示。在其中取 \(u=u_\sigma\)。反复使用乘法规则，可知 \(u^i\) 是 \(u_{\sigma^i}\) 的非零 \(K\) 倍数，特别地 \(u^n=a\in K^\times\)。因为 \(u\) 与其幂 \(a\) 交换，而 \(ua=\sigma(a)u\)，可消去 \(u\) 得 \(\sigma(a)=a\)，所以 \(a\in k^\times\)。将原基换成 \(1,u,\ldots,u^{n-1}\)，就是一次前述缩放基元的操作；它不改变余圈的上同调类，却将乘法误差化为
\[
u^i u^j=
\begin{cases}
u^{i+j},&i+j<n,\\
a\,u^{i+j-n},&i+j\ge n,
\end{cases}
\qquad(0\le i,j<n).
\]
因此相应余圈为
\[
c_a(\sigma^i,\sigma^j)=a^{\lfloor(i+j)/n\rfloor},
\]
且 \(B_c\cong(K/k,\sigma,a)\)。这证明参数给出从 \(k^\times\) 到 \(\operatorname{Br}(K/k)\) 的满射。进位余圈满足 \(c_a c_b=c_{ab}\)，由\cref{b2:thm:relative-brauer-cocycles}，该参数映射是群同态。

若再用 \(u'=bu\)、\(b\in K^\times\)，则
\[
(bu)^n=b\,\sigma(b)\cdots\sigma^{n-1}(b)\,u^n
=\operatorname{N}_{K/k}(b)a.
\]
因而乘以一个范数正对应再次改变这组基元。由\cref{b2:thm:cyclic-algebra-norm-splitting}，参数映射的核恰为 \(\operatorname{N}_{K/k}(K^\times)\)。商群同构定理遂给出所述范数商，连同\cref{b2:thm:relative-brauer-cocycles}便得到全部同构。
\end{proof}

交叉积的基元除了记录乘法，也使约化迹可以直接计算。左乘 \(u_h\) 将指标 \(g\) 移到 \(hg\)，因此只有单位指标的项会贡献矩阵的对角元；计算时仍须把左 \(K\) 系数转换成右 \(K\) 坐标。

\begin{proposition}\label{b2:prop:crossed-product-reduced-trace}
设 \(K/k\) 为有限 Galois 扩张，\(G=\operatorname{Gal}(K/k)\)，\(c:G\times G\to K^\times\) 为归一化 \(2\)-余圈，\(B_c=\bigoplus_{g\in G}Ku_g\) 为相应交叉积代数。对任意 \(x=\sum_{g\in G}a_g u_g\)、\(a_g\in K\)，有
\[
\operatorname{Trd}_{B_c/k}(x)=\operatorname{Tr}_{K/k}(a_1)
=\sum_{g\in G}g(a_1).
\]
\end{proposition}
\begin{proof}
记 \(n=[K:k]\)，将 \(B_c\) 的右 \(K\) 空间记为 \(W\)。由 \(a u_g=u_g g^{-1}(a)\)，\((u_g)_{g\in G}\) 也是 \(W\) 的一组基。根据\cref{b2:prop:maximal-field-splits-csa}中的分裂同构，
\[
B_c\otimes_k K\longrightarrow
\operatorname{End}_K(W),\qquad
x\otimes b\longmapsto(y\mapsto xyb),
\]
其中目标作用在原代数的右 \(K\) 空间上，而来源是标量扩张代数。由于 \(\dim_KW=n\)，右端是 \(M_n(K)\) 的一个模型。在这个模型中，\(x\otimes1\) 是左乘 \(x\) 的算子。对 \(h,g\in G\)，
\[
(a_hu_h)u_g
=a_h c(h,g)u_{hg}
=u_{hg}(hg)^{-1}\bigl(a_h c(h,g)\bigr).
\]
因此在列 \(g\) 中，这一项的矩阵系数只出现在行 \(hg\)。若 \(h\ne1\)，行指标不等于列指标，故它不贡献对角元；若 \(h=1\)，由 \(c(1,g)=1\)，对角元恰为 \(g^{-1}(a_1)\)。左乘 \(x\) 的矩阵迹于是为
\[
\sum_{g\in G}g^{-1}(a_1)=\sum_{g\in G}g(a_1)
=\operatorname{Tr}_{K/k}(a_1).
\]
有限 Galois 扩张的域迹等于全部共轭之和，而\cref{b2:def:reduced-trace-norm}中的约化迹就是任意分裂模型里的矩阵迹，故得所述公式。
\end{proof}

第五卷第 5.3 节的 Brauer 诱导定理研究有限群特征标，与本章按张量积运算的 Brauer 群是不同对象。

% !TeX root = ../../Book2.tex
% 纲要标题：### 18.5 周期与指数
% 纲要要点（供目录核对）：
%
% * 周期整除指数，指数整除次数
% * 十六维中心除代数的周期二、指数四例子
%
% <!-- 短节：不设小节 -->
%
% ---

\section{周期与指数}
\label{b2:sec:1805}

指数是除代数代表的次数，也等于有限分裂域扩张次数的最小值；周期则是 Brauer 类的加法阶。前节的半线性矩阵把两种量放在同一计算中：取矩阵行列式，会使因子集的一个幂成为余边。

\begin{theorem}\label{b2:thm:period-divides-index}
设 \(k\) 为域，\(A\) 为中心单 \(k\) 代数，则其 Brauer 类 \([A]\in\operatorname{Br}(k)\) 具有有限阶，且
\[
\operatorname{per}(A)\mid\operatorname{ind}(A)\mid\deg A.
\]
\end{theorem}
\begin{proof}
取同一 Brauer 类中的除代数 \(D\)，令 \(d=\deg D=\operatorname{ind}(A)\)，再取一个有限 Galois 分裂域 \(K/k\)。由\cref{b2:prop:split-algebra-factor-set}，对 \(D_K\cong M_d(K)\) 选矩阵 \(C_g\)，得到
\[
C_g\cdot g(C_h)=c(g,h)C_{gh},\qquad [D]=-[B_c].
\]
取矩阵行列式，令 \(b_g=\det C_g\)，则
\[
c(g,h)^d=b_g\cdot g(b_h)b_{gh}^{-1}.
\]
因此 \(c^d\) 为余边，在 \(H^2(G,K^\times)\) 中为零类。\cref{b2:thm:relative-brauer-cocycles}给出 \(d[B_c]=0\)，所以 \(d[A]=d[D]=0\)。元素的最小正阶于是存在，且整除任何零化它的正整数，特别整除 \(d\)。若 \(A\cong M_r(D)\)，则 \(\deg A=rd\)，还得到第二个整除关系。证明中没有除以 \(d\)，故包括正特征。
\end{proof}

\begin{example}[周期为二、指数为四]\label{b2:exm:period-two-index-four}
在域 \(k=\mathbb R((x))((y))\) 上，代数
\[
D=\mathbb H\otimes_{\mathbb R}k\ \otimes_k\ (x,y)_k
\]
为除代数，且 \(\operatorname{ind}(D)=4\)、\(\operatorname{per}(D)=2\)。这里 \((x,y)_k\) 为四元数代数。
\end{example}
\begin{proof}
给出一个直接的除环模型。先取中心不定元 \(u\)，组成 Laurent 级数除环 \(F=\mathbb H((u))\)。非零级数可以提出其首项，再将剩余的 \(1+h\)、\(h\) 只有正次数，按 \(1-h+h^2-\cdots\) 求逆；每个固定次数的系数只涉及有限多项，所以这个形式逆有意义。

令 \(\tau\) 固定 \(\mathbb H\)，将 \(u\) 送到 \(-u\)。在所有下端有界的形式级数 \(\sum_m a_m v^m\)、\(a_m\in F\) 上规定
\[
(a v^r)(b v^s)=a\tau^r(b)v^{r+s}.
\]
固定次数的系数仍只涉及有限和；\(\tau\) 是自同构，所以三项乘积的两种括号给出相同系数。这便定义结合环 \(E=F((v;\tau))\)。非零首项 \(a_m v^m\) 可逆，其逆为 \(\tau^{-m}(a_m^{-1})v^{-m}\)；提出首项后，再用同样的正次数几何级数构造逆。因此 \(E\) 是除环。

在 \(E\) 中令 \(x=u^2,y=v^2\)。二者都在中心，且 \(F\) 中按 \(u\) 的奇偶次数及四元数基逐项分组，得到左 \(\mathbb R((x))\) 空间的直和
\[
F=\bigoplus_{\epsilon=0}^{1}
\ \bigoplus_{h\in\{1,\mathbf i,\mathbf j,\mathbf k\}}
\mathbb R((x))\,h u^\epsilon.
\]
由于 \(\tau^2=\operatorname{id}\)，\(y=v^2\) 与 \(F\) 交换。再按 \(v\) 的奇偶次数分组，就有
\[
E=F((y))\oplus F((y))v
=\bigoplus_{\delta=0}^{1}
\ \bigoplus_{\epsilon=0}^{1}
\ \bigoplus_{h\in\{1,\mathbf i,\mathbf j,\mathbf k\}}
k\,h u^\epsilon v^\delta,
\qquad k=\mathbb R((x))((y)).
\]
两次形式 Laurent 展开的系数唯一，且 \(1,\mathbf i,\mathbf j,\mathbf k\) 在 \(\mathbb H\) 中线性无关，因此所列十六个元素确为 \(E\) 的 \(k\) 基。

中心也可逐项确定。若 \(z=\sum_m a_m v^m\) 与全部四元数常数交换，因 \(v\) 与这些常数交换，Laurent 展开的唯一性迫使每个 \(a_m\) 与 \(\mathbb H\) 交换。再逐项展开 \(a_m\) 的 \(u\) 系数，得到 \(a_m\in\mathbb R((u))\)。若 \(z\) 还与 \(u\) 交换，由 \(v^m u=(-1)^muv^m\) 及特征为零，所有奇数 \(m\) 的系数都为零。最后与 \(v\) 交换，使每个余下的 \(a_{2r}\) 被 \(\tau\) 固定；其 \(u\) 展开只能含偶数次幂，故 \(a_{2r}\in\mathbb R((x))\)。因此 \(Z(E)=\mathbb R((x))((y))=k\)。

四元数常数与 \(u,v\) 交换，而 \(u^2=x,v^2=y,vu=-uv\)。这些关系及所列基给出 \(E\cong\mathbb H_k\otimes_k(x,y)_k=D\)。因此 \(D\) 自身为十六维中心除代数，指数为四。由\eqref{b2:eq:quaternion-conjugation}，每个四元数代数都通过共轭同构于自身的反代数，其 Brauer 类的二倍为零，所以 \(2[D]=0\)。\(D\) 是不等于 \(k\) 的除代数，Brauer 类非零，故周期恰为二。
\end{proof}


\begin{chaptersummary}
Brauer 类由唯一的中心除代数代表确定，张量积给出加法，反代数给出负元。次数是代数在分裂域上成为全矩阵代数后的矩阵大小；指数记录除代数代表的次数，也是所有有限分裂域扩张次数中的最小值。周期是类的群阶，整除指数，但两者不必相等。

\ProseParagraphBreak
循环代数把自同构与参数转成扭曲乘法，参数是范数恰好等价于分裂。约化迹和约化范数在任意分裂模型中都给出相同的底域元素，与标量扩张相容；它们仍是具体代数中的元素数据，稳定加矩阵块会使 \(a\otimes I_r\) 的约化迹成为 \(r\operatorname{Trd}(a)\)，范数成为 \(\operatorname{Nrd}(a)^r\)。正则作用的迹和行列式还带有代数次数的倍数或幂次，不能不加区别地替代约化运算。第一类对合保持约化特征多项式，并与约化迹合成一个非退化对称型；这个迹型还会反映所选对合的信息。

\ProseParagraphBreak
交叉积的结合律正对应\(2\)-余圈方程，改变基元对应乘一个余边。对有限 Galois 扩张 \(K/k\)，相对 Brauer 群 \(\operatorname{Br}(K/k)\) 由这些余圈类描述；半线性矩阵的行列式进一步给出周期整除指数。周期二、指数四的例子说明，即使张量平方已经分裂，一个二次扩张仍未必足以分裂原代数。
\end{chaptersummary}

\BookExercises

\begin{exercise}\label{b2:ex:18-explicit-stabilization}
令 \(A=M_2(\mathbb H)\)、\(B=M_3(\mathbb H)\)。求使 \(M_r(A)\cong M_s(B)\) 的全部正整数对 \((r,s)\)，并分别列出 \(A,B\) 的次数、指数、周期。它们是否同构？
\end{exercise}
\begin{solution}
两边分别为 \(M_{2r}(\mathbb H)\)、\(M_{3s}(\mathbb H)\)。矩阵块大小相等时同构，反向维数相等迫使 \(2r=3s\)。所以全部解为 \((r,s)=(3t,2t)\)、\(t\ge1\)。两代数的次数分别为四、六，指数均为二，周期也均为二，因为除代数代表都是 \(\mathbb H\)，而四元数共轭表明其非零 Brauer 类的二倍为零。原代数实维数分别为十六、三十六，所以并不同构。
\end{solution}

\begin{exercise}\label{b2:ex:18-tensor-index-period-bounds}
设 \(k\) 为域，\(A,B\) 为中心单 \(k\) 代数。证明
\[
\operatorname{ind}(A\otimes_kB)\mid\operatorname{ind}(A)\operatorname{ind}(B),
\]
\[
\operatorname{per}(A\otimes_kB)\mid
\operatorname{lcm}\bigl(\operatorname{per}(A),\operatorname{per}(B)\bigr).
\]
用一个例子说明两项整除都可以严格。
\end{exercise}
\begin{solution}
取除代数代表 \(D,E\)，次数分别为 \(d,e\)。\(A\otimes B\) 与 \(D\otimes E\) 同类，后者次数为 \(de\)，其指数整除次数，得到第一式。记两个周期为 \(m,n\)，则 \(\operatorname{lcm}(m,n)\) 在 Brauer 群中同时零化 \([A]\) 和 \([B]\)，也零化它们的和，所以该和的周期整除此数。

取 \(A=B=\mathbb H\)。由四元数共轭及\cref{b2:thm:csa-enveloping-isomorphism}，\(\mathbb H\otimes_{\mathbb R}\mathbb H\cong M_4(\mathbb R)\)，其指数、周期均为一，而两式右侧分别为四和二。
\end{solution}

\begin{exercise}\label{b2:ex:18-index-after-extension}
设 \(k\) 为域，\(A\) 为中心单 \(k\) 代数，\(\operatorname{ind}(A)=12\)。若 \(L/k\) 为次数十八的有限扩张，列出\cref{b2:prop:index-after-base-change}所允许的 \(\operatorname{ind}(A_L)\) 值；这些约束是否已经保证每个值都能实现？若另一个有限扩张 \(F/k\) 使指数降为三，其次数须满足什么整除条件？再说明次数五的扩张能否改变原指数。
\end{exercise}
\begin{solution}
记 \(e=\operatorname{ind}(A_L)\)、\(s=12/e\)。正文命题要求 \(s\) 同时整除十二与十八，所以 \(s\in\{1,2,3,6\}\)，相应 \(e\in\{12,6,4,2\}\)。这是必要条件给出的候选值；它没有构造具有指定指数的扩域，不能单凭这些整数关系断言各候选值都能实现。

若指数降为三，则 \(12/3=4\) 必整除 \([F:k]\)，因此例如次数六的扩张不能使指数降为三。次数五与十二互素，正文命题使扩域后的指数仍为十二。
\end{solution}

\begin{exercise}\label{b2:ex:18-cyclic-splitting-matrices}
设 \(C=(K/k,\sigma,a)\)，\([K:k]=n\ge2\)。在原代数 \(C\) 的右 \(K\) 基 \(1,u,\ldots,u^{n-1}\) 中写出左乘 \(b\in K\) 与左乘 \(u\) 的矩阵。由此求 \(u\) 的约化特征多项式、约化迹和约化范数。
\end{exercise}
\begin{solution}
由 \(bu^j=u^j\sigma^{-j}(b)\)，左乘 \(b\) 的矩阵为
\[
D_b=\operatorname{diag}\bigl(b,\sigma^{-1}(b),\ldots,\sigma^{-(n-1)}(b)\bigr).
\]
左乘 \(u\) 把前面的基向量依次移到后一个，把最后一个送到 \(a\)，所以矩阵为 \(T^n-a\) 的友矩阵：其下次对角元为一，右上角为 \(a\)，其余为零。这些矩阵正来自\cref{b2:thm:cyclic-algebra-central-simple}证明中的分裂模型。

由\cref{b2:lem:companion-characteristic}，约化特征多项式是 \(T^n-a\)。因此 \(\operatorname{Trd}(u)=0\)、\(\operatorname{Nrd}(u)=(-1)^{n-1}a\)。这个符号来自特征多项式常数项的约定；尤其次数二时，约化范数并不等于参数 \(a\) 本身。
\end{solution}

\begin{exercise}\label{b2:ex:18-f9-cyclic-matrices}
取 \(K=\mathbb F_9=\mathbb F_3[i]\)、\(i^2=2\)，\(\sigma(z)=z^3\)。构造 \((K/\mathbb F_3,\sigma,2)\cong M_2(\mathbb F_3)\) 的显式同构，写出 \(i,u\) 的矩阵，并核对 \(u\) 的约化迹、范数。
\end{exercise}
\begin{solution}
\(\sigma(i)=2i=-i\)，而 \(b=1+i\) 满足
\[
\operatorname{N}_{K/\mathbb F_3}(b)=(1+i)(1-i)=1-2=2.
\]
在基 \(1,i\) 下，让 \(i\) 左乘，让 \(u\) 按 \(z\mapsto b\cdot\sigma(z)\) 作用，得到
\[
I=\begin{pmatrix}0&2\\1&0\end{pmatrix},\qquad
U=\begin{pmatrix}1&1\\1&2\end{pmatrix}.
\]
直接相乘有 \(I^2=U^2=2I_2\)、\(UI=-IU\)，所以它们满足循环代数的全部关系。所得含幺同态由来源单性为单射，两边都是四维，故为同构。\(\operatorname{tr}U=0\)、\(\det U=1=-2\)，与上一题的迹范数公式相符。
\end{solution}

\begin{exercise}\label{b2:ex:18-cyclic-generator-opposite}
在 \(C=(K/k,\sigma,a)\) 中，若整数 \(r\) 与 \(n=[K:k]\) 互素，证明改用 \(v=u^r\) 给出
\[
C\cong(K/k,\sigma^r,a^r).
\]
再证明 \(C^{\mathrm{op}}\cong(K/k,\sigma^{-1},a)\cong(K/k,\sigma,a^{-1})\)。
\end{exercise}
\begin{solution}
\(u\) 可逆，所以允许负整数幂。\(v\) 满足 \(vb=\sigma^r(b)v\)、\(v^n=a^r\)。取整数 \(s,t\) 使 \(rs+nt=1\)，则 \(u=v^s a^t\)，故 \(K,v\) 生成整个 \(C\)。两边维数相同，关系给出的满同态为同构。

在反代数中，使用原来的 \(u\) 有 \(u\mathbin{\cdot_{\mathrm{op}}}b=bu=u\sigma^{-1}(b)=\sigma^{-1}(b)\mathbin{\cdot_{\mathrm{op}}}u\)，且其 \(n\) 次幂仍为 \(a\)，得到第一个同构。再在该循环代数中把生成元取逆，按前一部分的 \(r=-1\) 得到第二个同构。
\end{solution}

\begin{exercise}\label{b2:ex:18-rational-norm-brauer-classes}
设 \(K=\mathbb Q(i)\)、\(\sigma(i)=-i\)。证明 \(C=(K/\mathbb Q,\sigma,3)\) 是除代数，确定其指数与周期，并证明参数改为 \(27\) 后代数同构。再把它与 \(H=(-1,-1)_{\mathbb Q}\) 比较：两者指数、周期相同，是否 Brauer 等价？
\end{exercise}
\begin{solution}
若三是范数，就有有理数 \(x,y\) 满足 \(x^2+y^2=3\)。清分母并约去公因子，得到互素整数 \(r,s,t\)、\(t\ne0\)，满足 \(r^2+s^2=3t^2\)。模三的平方只有零和一，所以 \(3\mid r,s\)。代回后又得 \(3\mid t\)，矛盾。因此由\cref{b2:thm:cyclic-algebra-norm-splitting}，\(C\) 不分裂；它是次数二的四元数代数，故由\cref{b2:prop:quaternion-division-or-split}为除代数，指数为二。共轭给出与反代数的同构，非零 Brauer 类的周期因而恰为二。

把生成元 \(u\) 换成 \(3u\)，其平方为 \(27\)，且与 \(K\) 的交换关系不变，所以得到所述代数同构。

\(H\) 的范数是四个有理数平方之和，非零向量的范数不为零，所以它也为除代数，指数、周期均为二。但是扩张到 \(\mathbb R\) 后，\(C\) 因正参数三是复数范数而分裂，\(H\) 则成为不分裂的 \(\mathbb H\)。若原来的 Brauer 类相同，扩域后的类也应相同，矛盾。因此相同的指数与周期不足以确定 Brauer 类。
\end{solution}

\begin{exercise}\label{b2:ex:18-reduced-trace-pairing}
设 \(k\) 为特征 \(p>0\) 的域，\(A=M_p(k)\)，\(\beta(X,Y)=\operatorname{tr}(XY)\)。求矩阵单位基的 \(\beta\) 对偶基，并求迹零子空间 \(H=\{X\in A:\operatorname{tr}X=0\}\) 上限制型的根空间。解释原配对非退化、\(\operatorname{tr}(I_p)=0\) 与限制型退化为何可以同时成立。
\end{exercise}
\begin{solution}
矩阵单位满足 \(\beta(E_{ij},E_{ab})=\delta_{ja}\delta_{ib}\)，所以 \(E_{ij}\) 的对偶向量为 \(E_{ji}\)，原配对非退化。

若 \(X\in A\) 与全部 \(H\) 配对为零，先取 \(E_{ij}\in H\)、\(i\ne j\)，得到 \(X_{ji}=0\)，故 \(X\) 对角。再取 \(E_{ii}-E_{11}\in H\)，得到 \(X_{ii}=X_{11}\)。因此 \(H^\perp=kI_p\)；反向包含由 \(\beta(cI_p,Y)=c\operatorname{tr}Y\) 成立。

特征 \(p\) 使 \(\operatorname{tr}(I_p)=p=0\)，所以 \(kI_p\subseteq H\)。限制型的根为 \(H\cap H^\perp=kI_p\)，维数一。原迹映射仍满射，因为 \(\operatorname{tr}(E_{11})=1\)；零迹单位元只是与这个迹零子空间正交，不是与全部 \(A\) 正交。
\end{solution}

\begin{exercise}\label{b2:ex:18-quaternion-block-reduced-polynomial}
令 \(A=M_2(\mathbb H)\)、\(a=\operatorname{diag}(\mathbf i,1+\mathbf j)\)。求 \(a\) 的约化特征多项式、约化迹与范数，并求左乘 \(a\) 在底层实空间上的迹和行列式。
\end{exercise}
\begin{solution}
扩张到 \(\mathbb C\) 后，两对角块各变为二阶矩阵。\(\mathbf i\) 的约化特征多项式为 \(T^2+1\)；\(1+\mathbf j\) 的约化迹为二、范数为二，故其多项式为 \(T^2-2T+2\)。因此
\[
\chi^{\mathrm{red}}_{A,a}(T)
=(T^2+1)(T^2-2T+2)
=T^4-2T^3+3T^2-2T+2.
\]
得到 \(\operatorname{Trd}(a)=2\)、\(\operatorname{Nrd}(a)=2\)。\(\deg A=4\)，底层实维数为十六，由\cref{b2:prop:regular-reduced-characteristic}，左乘算子的迹为八、行列式为 \(2^4=16\)。
\end{solution}

\begin{exercise}\label{b2:ex:18-reduced-tensor-formulas}
设 \(k\) 为域，\(A,B\) 为中心单 \(k\) 代数，\(\deg A=m\)、\(\deg B=n\)。对 \(a\in A,b\in B\) 证明
\[
\operatorname{Trd}_{A\otimes B}(a\otimes b)
=\operatorname{Trd}_A(a)\operatorname{Trd}_B(b),
\]
\[
\operatorname{Nrd}_{A\otimes B}(a\otimes b)
=\operatorname{Nrd}_A(a)^n\operatorname{Nrd}_B(b)^m.
\]
并说明稳定扩张 \(a\mapsto a\otimes I_r\) 如何改变这两个值。
\end{exercise}
\begin{solution}
在共同的代数闭分裂域中，两边成为矩阵的 Kronecker 积\footnote{克罗内克（Leopold Kronecker，1823-1891），德国数学家，研究数论、代数与矩阵。}。其对角项逐对相乘，迹为两个迹之积。分别把两个矩阵化为上三角形；可以使用\cref{b2:thm:jordan-canonical-form}，不要求可对角化。若对角元分别为 \(\lambda_i,\mu_j\)，张量矩阵的对角元为全部 \(\lambda_i\mu_j\)，其积为 \((\prod_i\lambda_i)^n(\prod_j\mu_j)^m\)，得到范数公式。由\cref{b2:prop:reduced-base-change}，这些等式下降到 \(k\)。

取 \(B=M_r(k),b=I_r\)，得到新约化迹为 \(r\operatorname{Trd}_A(a)\)，新约化范数为 \(\operatorname{Nrd}_A(a)^r\)。因此 Brauer 类忽略矩阵大小，但具体元素的约化运算仍记录代数次数。
\end{solution}

\begin{exercise}\label{b2:ex:18-opposite-reduced-polynomial}
设 \(k\) 为域，\(A\) 为中心单 \(k\) 代数，将 \(a\in A\) 对应到反代数中的元素 \(a^{\mathrm{op}}\)。证明其约化特征多项式、迹和范数均不变，并解释这与 \([A^{\mathrm{op}}]=-[A]\) 不矛盾。
\end{exercise}
\begin{solution}
若 \(\rho:A_L\to M_n(L)\) 是分裂同构，则 \(a^{\mathrm{op}}\mapsto\rho(a)^{\mathsf T}\) 是反代数的分裂同构，因为转置也反转乘法。矩阵与其转置有相同的特征多项式，所以约化特征多项式及两个指定系数不变。Brauer 群的负元通过整个代数的张量积刻画，并不是把每个元素的约化迹或范数取相反数；这两种运算的对象不同。
\end{solution}

\begin{exercise}\label{b2:ex:18-inseparable-field-trace}
设 \(k\) 为特征 \(p>0\) 的域，\(E=k(\alpha)\)、\(\alpha^p\in k\)、\([E:k]=p\)。通过正则作用 \(a\mapsto m_a\) 将 \(E\) 嵌入 \(A=\operatorname{End}_k(E)\)。在基 \(1,\alpha,\ldots,\alpha^{p-1}\) 下，令 \(T_j\) 将 \(\alpha^j\) 送到一、其余基向量送到零，\(0\le j<p\)。计算 \(\operatorname{Trd}(m_aT_j)\)，并求 \(E\) 对约化迹配对的正交空间维数。比较 \(p=2\) 与 \(p>2\) 时 \(E\subseteq E^\perp\) 是否相等。
\end{exercise}
\begin{solution}
写 \(a=\sum_{r=0}^{p-1}c_r\alpha^r\)。算子 \(m_aT_j\) 只有第 \(j\) 列可能非零，该列正是 \(a\) 的坐标列，所以其迹等于 \(c_j\)。\(A\) 已为分裂矩阵代数，约化迹就是这个矩阵迹，故
\[
\operatorname{Trd}(m_aT_j)=c_j.
\]
这些算子分别读取 \(E\) 的全部坐标，说明 \(A\to E^*\)、\(T\mapsto(a\mapsto\operatorname{Trd}(m_aT))\) 满射。其核为 \(E^\perp\)，因此 \(\dim_kE^\perp=p^2-p\)。

正文例子的幂基计算在这里仍适用：乘 \(\alpha^i\)、\(1\le i<p\)，将每个基向量的指数移到不同的模 \(p\) 余数类，没有对角项；乘标量 \(c\in k\) 的迹则为 \(pc=0\)。因此 \(E\) 上全部乘法算子的迹均为零，\(E\) 上的限制配对也恒为零，故 \(E\subseteq E^\perp\)。当 \(p=2\) 时两者维数都是二，所以相等；当 \(p>2\) 时 \(p^2-p>p\)，所以包含严格。尽管 \(E\) 内部无法用迹检测元素，\(T_j\) 这些外部算子却能恢复每一个坐标。
\end{solution}

\begin{exercise}\label{b2:ex:18-norm-nonlinearity-nilpotents}
设 \(k\) 为域，在 \(M_2(k)\) 中令 \(X_t=\begin{psmallmatrix}t&1\\0&0\end{psmallmatrix}\)、\(t\in k\)。求其约化迹、范数，计算 \(\operatorname{Nrd}(X_t+E_{22})-\operatorname{Nrd}(X_t)-\operatorname{Nrd}(E_{22})\)，并判断哪些 \(X_t\) 幂零。再证明任意中心单代数中的幂零元素约化迹、范数均为零，并给出两者同时为零却不幂零的矩阵。
\end{exercise}
\begin{solution}
矩阵代数的约化运算为通常运算，所以 \(\operatorname{Trd}(X_t)=t\)、\(\operatorname{Nrd}(X_t)=0\)。加上 \(E_{22}\) 后行列式为 \(t\)，而 \(\operatorname{Nrd}(E_{22})=0\)，所求差正是 \(t\)。例如 \(t=1\) 就显示范数不保持加法。

\(X_t^2=tX_t\)，归纳得 \(X_t^r=t^{r-1}X_t\)、\(r\ge1\)。因此 \(t=0\) 时平方为零，\(t\ne0\) 时每个正次幂均非零，恰好只有 \(X_0\) 幂零。

若中心单代数中 \(a^r=0\)，扩张到代数闭分裂域后所得矩阵的全部特征值都为零，所以约化特征多项式为 \(T^n\)，约化迹与范数均为零。反向失败：任意域上的 \(\operatorname{diag}(1,-1,0)\) 迹与行列式均为零，但每个正次幂仍有首个对角元一，故不幂零。
\end{solution}

\begin{exercise}\label{b2:ex:18-factor-set-changes}
设 \(K/k\) 为次数 \(n\ge2\) 的循环 Galois 扩张，生成元为 \(\sigma\)，\(a\in k^\times\)，\(A=(K/k,\sigma,a)\)。沿用\cref{b2:ex:18-cyclic-splitting-matrices}的右 \(K\) 基分裂模型，记左乘 \(b\in K\) 的矩阵为 \(D_b\)，左乘 \(u\) 的矩阵为 \(U\)。证明自然 Galois 作用可取实现矩阵 \(C_{\sigma^i}=U^{-i}\)、\(0\le i<n\)，计算其因子集，并说明它如何检验\cref{b2:prop:split-algebra-factor-set}中的负号。
\end{exercise}
\begin{solution}
这里 \(D_b\) 的第 \(j\) 个对角元为 \(\sigma^{-j}(b)\)、\(0\le j<n\)，而 \(U\) 将基向量循环移到下一位，越过末位时乘以 \(a\)。所以
\[
U D_bU^{-1}=D_{\sigma(b)},\qquad
\sigma(D_b)=D_{\sigma(b)},\qquad \sigma(U)=U.
\]
因此 \(X\mapsto U^{-1}\sigma(X)U\) 固定原代数的全部 \(D_b,U\)，并在标量上施加 \(\sigma\)，正是自然半线性作用。其第 \(i\) 次迭代由 \(U^{-i}\) 实现。

所有 \(C_{\sigma^i}\) 的条目都在 \(k\) 中。若 \(i+j=qn+r\)、\(0\le r<n\)，则 \(U^n=aI_n\) 给出
\[
C_{\sigma^i}\sigma^i(C_{\sigma^j})
=U^{-(i+j)}=a^{-q}U^{-r}.
\]
故因子集为 \(c(\sigma^i,\sigma^j)=a^{-\lfloor(i+j)/n\rfloor}\)，它是参数 \(a^{-1}\) 的进位余圈，而不是参数 \(a\) 的进位余圈。相应交叉积为 \((K/k,\sigma,a^{-1})\)，所以 \([A]=-[B_c]\) 恰好把该逆参数还原为原来的 Brauer 类。
\end{solution}

\begin{exercise}\label{b2:ex:18-cyclic-cohomology}
设 \(K/k\) 为次数 \(n=rm\) 的循环 Galois 扩张，\(r,m\ge2\)，生成元为 \(\sigma\)。令 \(H=\langle\sigma^r\rangle\)、\(F=K^H\)。将参数 \(a\in k^\times\) 的进位余圈限制到 \(H\)，求所得参数。证明自然包含 \(k^\times\hookrightarrow F^\times\) 诱导范数商群之间的同态，并说明它对应余圈的限制。
\end{exercise}
\begin{solution}
\(H\) 的阶为 \(m\)，且 \(K/F\) 的生成自同构为 \(\sigma^r\)。对 \(0\le i,j<m\)，限制后的余圈满足
\[
c_a(\sigma^{ri},\sigma^{rj})
=a^{\lfloor r(i+j)/(rm)\rfloor}
=a^{\lfloor(i+j)/m\rfloor},
\]
仍是参数 \(a\) 的进位余圈，只是循环长度由 \(n\) 变成 \(m\)。

若 \(a\) 在原范数陪集中乘上 \(N_{K/k}(b)\)，令 \(z=\prod_{j=0}^{r-1}\sigma^j(b)\)，则
\[
N_{K/F}(z)
=\prod_{i=0}^{m-1}\prod_{j=0}^{r-1}\sigma^{ri+j}(b)
=N_{K/k}(b).
\]
因此原范数子群包含于 \(N_{K/F}(K^\times)\)，得到良定义群同态
\[
k^\times/N_{K/k}(K^\times)
\longrightarrow F^\times/N_{K/F}(K^\times),
\qquad [a]\longmapsto[a].
\]
余圈和余边限制到子群后仍满足同一方程；所算的进位代表说明，在正文的循环范数商识别下，这个同态正对应限制 \(H^2(\langle\sigma\rangle,K^\times)\to H^2(H,K^\times)\)。
\end{solution}

\begin{exercise}\label{b2:ex:18-real-brauer-group}
利用\cref{b2:cor:cyclic-relative-brauer}及复数范数，计算 \(\operatorname{Br}(\mathbb R)\)，并据此确定全部有限维实中心除代数的同构类型。不使用实除代数分类定理。
\end{exercise}
\begin{solution}
任意实中心单代数扩域到 \(\mathbb C\) 后仍中心单，而代数闭域上的中心单代数为全矩阵代数，见\cref{b2:thm:csa-base-change}及\cref{b2:prop:csa-division-description}。因此所有实 Brauer 类都由 \(\mathbb C\) 分裂，即 \(\operatorname{Br}(\mathbb R)=\operatorname{Br}(\mathbb C/\mathbb R)\)。由\cref{b2:thm:relative-brauer-cocycles}和\cref{b2:cor:cyclic-relative-brauer}，
\[
\operatorname{Br}(\mathbb R)
\cong\mathbb R^\times/\operatorname{N}_{\mathbb C/\mathbb R}(\mathbb C^\times)
=\mathbb R^\times/\mathbb R_{>0}
\cong\mathbb Z/2\mathbb Z.
\]
正参数给零类，负参数给 \(\mathbb H\) 的类。由\cref{b2:prop:brauer-division-representative}的除代数代表唯一性，两个类的中心除代数代表分别为 \(\mathbb R\) 和 \(\mathbb H\)。因此全部实中心单代数是某个 \(M_r(\mathbb R)\) 或 \(M_r(\mathbb H)\)，结论来自 Brauer 对应及已知分裂性。
\end{solution}

\begin{exercise}\label{b2:ex:18-quaternion-factor-set}
在 \(\mathbb H\otimes_{\mathbb R}\mathbb C\cong M_2(\mathbb C)\) 中取
\[
\mathbf i\longmapsto I=\begin{pmatrix}i&0\\0&-i\end{pmatrix},\qquad
\mathbf j\longmapsto J=\begin{pmatrix}0&-1\\1&0\end{pmatrix}.
\]
令 \(s\) 为复共轭。证明自然半线性作用为 \(\alpha_s(X)=J\overline XJ^{-1}\)，并计算因子集的 \(c(s,s)\)。
\end{exercise}
\begin{solution}
逐条目共轭将 \(I\) 送到 \(-I\)，将 \(J\) 固定；而 \(JI=-IJ\)，故再与 \(J\) 共轭以后，\(I,J\) 都被固定。该半线性作用在标量上是复共轭，因而正是固定原四元数代数的自然 Galois 作用。

取 \(C_s=J,C_1=I_2\)，则 \(C_s\overline{C_s}=J^2=-I_2\)，所以 \(c(s,s)=-1\)。相应交叉积是参数负一的循环代数，即 \(\mathbb H\)。在这个二阶 Brauer 类上，\([\mathbb H]=-[\mathbb H]\)，因此一般公式中的负号在该例上不改变结果。
\end{solution}

\begin{exercise}\label{b2:ex:18-crossed-product-reduced-trace}
设 \(K/k\) 为有限 Galois 扩张，\(G=\operatorname{Gal}(K/k)\)，\(c\) 为归一化 \(2\)-余圈，\(B_c=(K/k,G,c)\)。计算约化迹配对在齐次分量 \(Ku_g\times Ku_h\) 上的值。给定 \(K/k\) 的基 \((a_i)\) 及其域迹对偶基 \((a_i^*)\)，为 \(B_c\) 的基 \((a_i u_g)_{i,g}\) 写出约化迹对偶基。
\end{exercise}
\begin{solution}
由交叉积乘法与\cref{b2:prop:crossed-product-reduced-trace}，
\[
\operatorname{Trd}\bigl((a u_g)(b u_h)\bigr)=
\begin{cases}
0,&gh\ne1,\\
\operatorname{Tr}_{K/k}\bigl(a\,g(b)c(g,g^{-1})\bigr),&h=g^{-1}.
\end{cases}
\]
因此分量 \(Ku_g\) 只与逆指标分量 \(Ku_{g^{-1}}\) 配对。将
\[
w_{i,g}=g^{-1}\bigl(a_i^*c(g,g^{-1})^{-1}\bigr)u_{g^{-1}}
\]
代入，得到 \(\operatorname{Trd}((a_j u_h)w_{i,g})=\delta_{h,g}\operatorname{Tr}_{K/k}(a_j a_i^*)=\delta_{h,g}\delta_{j,i}\)。故 \((w_{i,g})\) 是所求对偶基。域迹对偶基的存在来自 \(K/k\) 可分时域迹配对非退化，见第一卷定理18.3.3及推论18.3.4。
\end{solution}

\begin{exercise}\label{b2:ex:18-period-index-minimal-splitting}
沿用\cref{b2:cor:period-index-biquadratic-tower}中的 \(k,D,i,s\)。设 \(F/k\) 为奇数次数的有限扩张，置于同一个代数闭包中。求 \(D_F\) 的指数与周期，并求 \(D_{F(i)}\)、\(D_{F(i,s)}\) 的指数和两级扩张次数。证明 \(F(i,s)/F\) 仍是最小有限分裂扩张。
\end{exercise}
\begin{solution}
令 \(e=\operatorname{ind}(D_F)\)。\cref{b2:prop:index-after-base-change}使 \(4/e\) 同时整除四与奇数 \([F:k]\)，所以 \(e=4\)。原来的 \(2[D]=0\) 经标量扩张仍成立，而指数四排除零类，故周期仍为二。

\(-1\) 不可能已经在 \(F\) 中有平方根，否则 \(k(i)\subseteq F\) 会使二整除 \([F:k]\)。因此 \([F(i):F]=2\)。这个扩张使四元数因子 \(\mathbb H_F\) 分裂，余类由次数二的另一个因子代表，指数至多为二；但不能为一，否则指数四会整除二次分裂次数，所以恰为二。

再添入 \(s\) 使第二因子也分裂，\(D_{F(i,s)}\) 的指数为一。\([F(i,s):F(i)]\le2\)，且不能为一，否则前一步已经分裂，故等于二。合成域的次数为四，指数四又排除次数更小的有限分裂域，因此它仍为最小分裂扩张。
\end{solution}

\begin{exercise}\label{b2:ex:18-finite-field-cocycles}
设 \(q\) 为素数幂、\(n\ge1\)，\(K=\mathbb F_{q^n}\)、\(k=\mathbb F_q\)。用循环群的范数商证明 \(H^2\bigl(\operatorname{Gal}(K/k),K^\times\bigr)=0\)，并说明如何把一个进位余圈的参数改成一。
\end{exercise}
\begin{solution}
\(n=1\) 时群平凡，归一化余圈只有平凡类。设 \(n\ge2\)。第一卷推论17.2.1和定理17.2.2分别说明 \(K^\times\) 循环、Galois 群由 \(z\mapsto z^q\) 生成。范数为
\[
z\longmapsto z^{1+q+\cdots+q^{n-1}}.
\]
若 \(z\) 为 \(K^\times\) 的生成元，其范数的阶为 \(q-1\)，所以范数满射到 \(k^\times\)。由\cref{b2:cor:cyclic-relative-brauer}，所求上同调群为零。

具体地，进位参数为 \(a\) 时，选 \(b\in K^\times\) 满足 \(\operatorname{N}_{K/k}(b)=a^{-1}\)，将循环生成元换成 \(bu\)，其 \(n\) 次幂便为一。再用它的各次幂作全部基向量，所有进位系数都成为一，所以余圈化为平凡余圈。
\end{solution}

\begin{exercise}\label{b2:ex:18-two-involution-trace-forms}
在 \(A=M_2(\mathbb R)\) 上分别取 \(\sigma_+(X)=X^{\mathsf T}\) 与 \(\sigma_J(X)=J^{-1}X^{\mathsf T}J\)，其中 \(J=\begin{psmallmatrix}0&1\\-1&0\end{psmallmatrix}\)。求两种迹型 \(t_\sigma(X,Y)=\operatorname{tr}\bigl(\sigma(X)Y\bigr)\) 的惯性，说明它们不是等距型。
\end{exercise}
\begin{solution}
写 \(X=\begin{psmallmatrix}a&b\\c&d\end{psmallmatrix}\)。第一种有 \(t_{\sigma_+}(X,X)=a^2+b^2+c^2+d^2\)，惯性为 \((4,0,0)\)。第二种直接计算为
\[
\sigma_J(X)=\begin{pmatrix}d&-b\\-c&a\end{pmatrix},\qquad
t_{\sigma_J}(X,X)=2(ad-bc).
\]
令 \(a=s+t,d=s-t,b=u+v,c=u-v\)，这是可逆实线性换元，所得型为 \(2s^2-2t^2-2u^2+2v^2\)，惯性为 \((2,2,0)\)。由实惯性定律，两型不等距。底层代数始终是同一个 \(M_2(\mathbb R)\)，差别来自对合的选择。
\end{solution}

\begin{exercise}\label{b2:ex:18-prime-index-period}
设 \(k\) 为域，\(p\) 为素数。证明指数为 \(p\) 的中心单 \(k\) 代数，其周期必为 \(p\)。由此给出次数为素数的中心单代数的两种可能。反过来，周期为素数是否迫使指数等于周期？
\end{exercise}
\begin{solution}
由\cref{b2:thm:period-divides-index}，周期整除 \(p\)，只能为一或 \(p\)。周期为一表示 Brauer 类为零，除代数代表就是底域，指数应为一，与题设矛盾。所以周期为 \(p\)。

次数为素数时，写 \(A=M_r(D)\)，有 \(p=r\deg D\)。若 \(\deg D=1\)，则 \(A=M_p(k)\) 分裂；否则 \(r=1\)，\(A\) 本身是次数 \(p\) 的除代数，指数、周期均为 \(p\)。逆向断言则被\cref{b2:exm:period-two-index-four}否定：其中周期二为素数，指数却为四。
\end{solution}
